\exerciseshdr{习题}

\exercisegroup

\begin{enumerate}
\def\labelenumi{\arabic{enumi})}
\tightlist
\item
  设 \((\mathrm{Q}_{i})_{i\in\mathrm{I}}\) 是一族箭图。用 S 表示族 \((\mathrm{Som}(\mathrm{Q}_{i}))\) 的和集，用 F 表示族 \((\mathrm{Fl}(\mathrm{Q}_{i}))\) 的和集，并用 o 与 t 表示由诸箭图 \(Q_{i}\) 的起点映射与终点映射诱导的从 F 到 S 的映射。
\end{enumerate}

\begin{enumerate}
\def\labelenumi{\alph{enumi})}
\item
  证明 \(\mathrm{Q} = (\mathrm{S}, \mathrm{F}, o, t)\) 是一个箭图。称它为族 \((\mathrm{Q}_{i})\) 的和箭图。
\item
  对 \(i \in I\)，存在唯一的箭图态射 \(j_{i}\)（由 \(Q_{i}\) 到 Q），使得映射 \(\operatorname{Som}(j_{i})\) 是从 \(\operatorname{Som}(Q_{i})\) 到 S 中的典范映射，且映射 \(\operatorname{Fl}(j_{i})\) 是从 \(\operatorname{Fl}(Q_{i})\) 到 S 中的典范映射。
\item
  设 K 是一个箭图，且对每个 \(i \in I\)，设 \(\varphi_{i}\) 是从 \(Q_{i}\) 到 K 的箭图态射。证明存在唯一的从 Q 到 K 的箭图态射 \(\varphi\)，使得 \(\varphi \circ j_{i} = \varphi_{i}\) 对所有 \(i \in I\) 成立。
\end{enumerate}

\begin{enumerate}
\def\labelenumi{\arabic{enumi})}
\setcounter{enumi}{1}
\tightlist
\item
  设 \(Q = (S, F, o, t)\) 是一个箭图。对每个 \(s \in S\)，称 \(d_{+}(s)\)（相应地，\(d_{-}(s)\)）为 Q 在 s 处的出次数（相应地，入次数），并把 \(d_{+}(s)\)（相应地，\(d_{-}(s)\)）定义为集合 \(\{f \in F \mid o(f) = s\}\)（相应地，\(\{f \in F \mid t(f) = s\}\)）的基数。若 \(d_{+}(s)\) 与 \(d_{-}(s)\) 对每个 \(s \in S\) 均有限，则称 Q 是局部有限的；若集合 F 与 S 均有限，则称 Q 是有限的。
\end{enumerate}

证明等式 \(\sum_{s\in S}d_{+}(s) = \sum_{s\in S}d_{-}(s) = \mathrm{Card}(\mathrm{F})\)

\begin{enumerate}
\def\labelenumi{\arabic{enumi})}
\setcounter{enumi}{2}
\tightlist
\item
  设 \(\mathrm{Q} = (\mathrm{S},\mathrm{F},o,t)\) 是一个局部有限箭图。我们定义矩阵 \(M_{\mathrm{Q}} = (m_{i,j})\)（类型为 (S,S)，元素在 \(\mathbf{Z}\) 中）如下：
\end{enumerate}

\[
 m _ {i, j} = \operatorname{Card} \left\{f \in \mathrm{F} \mid o (f) = i \text {and} t (f) = j \right\}.
\]

矩阵 \(M_{\mathrm{Q}}\) 称为 Q 的邻接矩阵。

\begin{enumerate}
\def\labelenumi{\alph{enumi})}
\item
  给定箭图 Q 与 Q'，给出关于它们的邻接矩阵的一个必要充分条件，使 Q 与 Q' 同构。
\item
  用邻接矩阵 \(M_{Q}\) 表出 Q 中起点与终点给定的长为 n 的道路的数目。
\item
  证明一个图的底箭图的邻接矩阵是对称的。
\end{enumerate}

\begin{enumerate}
\def\labelenumi{\arabic{enumi})}
\setcounter{enumi}{3}
\tightlist
\item
  设 \(Q = (S, F, o, t)\) 是一个箭图。用 I 表示族 \((S, F)\) 的和集；把 S 与 F 等同于 I 的子集。设 k 是一个交换环。把箭图 Q 的 k-代数称为商 \(A_{Q}\)（即集合 I 上的自由结合 k-代数对由元素 \(X_{s}^{2} - X_{s}\)（对 \(s \in S\)）、\(X_{u}X_{v} - X_{v}X_{u}\)（对 \(u, v \in S\)）、\(X_{f}X_{t(f)} - X_{f}\) 与 \(X_{o(f)}X_{f} - X_{f}\)（对 \(f \in F\)）所生成的双边理想所作的商）。对 \(i \in I\)，用 \(x_{i}\) 表示 \(X_{i}\) 在代数 \(A_{Q}\) 中的像。
\end{enumerate}

\begin{enumerate}
\def\labelenumi{\alph{enumi})}
\item
  证明有 \(\sum_{s\in S}x_s = 1\)
\item
  证明若 f 与 g 是 Q 中不可合成的箭，则有 \(x_{f}x_{g}=0\)。
\item
  对 Q 中每条道路 \(c = (a_{0}, f_{1}, \ldots, f_{n}, a_{n})\)，令 \(x_{c} = x_{a_{0}} x_{f_{1}} \ldots x_{f_{n}}\)。设 c 与 \(c'\) 是 Q 中的道路；证明有 \(x_{c} x_{c'} = x_{cc'}\)（若 c 与 \(c'\) 可合成），否则有 \(x_{c} = x_{c'}\)。
\item
  设 C 是箭图 Q 中道路的集合；证明族 \((x_{c})_{c\in C}\) 是 k-模 \(A_{Q}\) 的一组基。
\item
  设 \(J_{Q}\) 是 \(A_{Q}\) 中由形如 \(x_{f}\)（\(f \in F\)）的元素生成的双边理想。证明代数 \(A_{Q}/J_{Q}\) 同构于 \(k^{(S)}\)。
\item
  为使 k-代数 \(A_{Q}\) 是一个有限生成的 k-模，必要且充分的是 Q 的顶点集有限且 Q 中每个圈的长为零。此时理想 \(J_{Q}\) 是幂零的。
\end{enumerate}

在 §1 的其余习题中，我们沿用这些记号。当我们说到 \(A_{Q}\)-模（不作另外说明时），指的是右模。若 A 是一个环，A-模 M（不作另外说明时）总指右 A-模；若 a 是 A 的一个元素，用 \(a_{M}\) 或 \((a)_{M}\) 表示 M 的位似 \(x \mapsto xa\)。

\begin{enumerate}
\def\labelenumi{\arabic{enumi})}
\setcounter{enumi}{4}
\tightlist
\item
  \begin{enumerate}
  \def\labelenumii{\alph{enumii})}
  \tightlist
  \item
    证明四元组 \(Q^{\circ} = (S, F, t, o)\) 是一个箭图。称它为 Q 的反箭图。
  \end{enumerate}
\end{enumerate}

\begin{enumerate}
\def\labelenumi{\alph{enumi})}
\setcounter{enumi}{1}
\tightlist
\item
  代数 \(A_{Q^{\circ}}\) 同构于代数 \(A_{Q}\) 的反代数。
\end{enumerate}

\begin{enumerate}
\def\labelenumi{\arabic{enumi})}
\setcounter{enumi}{5}
\tightlist
\item
  \begin{enumerate}
  \def\labelenumii{\alph{enumii})}
  \tightlist
  \item
    设 \((M_s)_{s\in\operatorname{Som}(Q)}\) 是一族 \(k\)-模，\((u_f)_{f\in\operatorname{Fl}(Q)}\) 是一个族，其中对每个 \(f\in\operatorname{Fl}(Q)\)，\(u_f\) 是从 \(M_{o(f)}\) 到 \(M_{t(f)}\) 的线性映射。设 M 是族 \((M_s)\) 的直和；对 \(s\in S\)，用 \(j_s\) 表示 \(M_s\) 到 M 中的典范单射，用 \(p_s\) 表示 M 到 \(M_s\) 上的典范投影。证明在 M 上存在唯一的 \(A_Q\)-模结构，使得有 \((x_s)_M = j_s \circ p_s\) 且 \((x_f)_M = j_{t(f)} \circ u_f \circ p_{o(f)}\)（对所有 \(s \in \operatorname{Som}(Q)\) 及每个 \(f \in \operatorname{Fl}(Q)\)）。
  \end{enumerate}
\end{enumerate}

反之，每个 \(A_{Q}\)-模都具有此形式。

\begin{enumerate}
\def\labelenumi{\alph{enumi})}
\setcounter{enumi}{1}
\tightlist
\item
  若 \(\Lambda\) 是一个 \(k\)-模且 \(s \in \operatorname{Som}(\mathbf{Q})\)，用 \(\Lambda(s)\) 表示如下的 \(\mathrm{A}_{\mathrm{Q}}\)-模：它与族 \((\mathrm{M}_i)_{i \in \operatorname{Som}(\mathbf{Q})}\)（一个 \(k\)-模族，由 \(\mathrm{M}_i = \Lambda\)（若 \(i = s\)）、否则 \(\mathrm{M}_i = 0\) 给出）以及线性映射族 \((u_f)_{f \in \operatorname{Fl}(\mathbf{Q})}\)（其中 \(u_f = 0\) 对一切 \(f \in \operatorname{Fl}(\mathbf{Q})\)）相伴。
\end{enumerate}

设 \(\Lambda\) 与 \(\Lambda'\) 是非零 \(k\)-模。设 \(s\) 与 \(s'\) 是 Q 的顶点。证明 \(\Lambda(s)\) 同构于 \(\Lambda'(s')\) 当且仅当 \(\Lambda\) 与 \(\Lambda'\) 作为 \(k\)-模同构且 \(s = s'\)。

\begin{enumerate}
\def\labelenumi{\alph{enumi})}
\setcounter{enumi}{2}
\item
  若 \(\Lambda\) 是一个单 \(k\)-模，则 \(\Lambda(s)\) 都是单 \(\mathrm{A}_{\mathrm{Q}}\)-模（对每个顶点 \(s\)）。
\item
  进一步假设 Q 中每个圈的长为零。证明每个单 \(A_Q\)-模都具有形式 \(\Lambda(s)\)，其中 \(\Lambda\) 是一个单 \(k\)-模而 s 是 Q 的一个顶点。
\end{enumerate}

\begin{enumerate}
\def\labelenumi{\arabic{enumi})}
\setcounter{enumi}{6}
\tightlist
\item
  设 \(k\) 是一个代数闭域。对每个 \(s \in \operatorname{Som}(Q)\)，令 \(\mathrm{P}(s) = x_s\mathrm{A}_{\mathrm{Q}}\)。
\end{enumerate}

\begin{enumerate}
\def\labelenumi{\alph{enumi})}
\item
  证明 \(\mathrm{P}(s)\) 是一个投射且不可分解的 \(A_{Q}\)-模，且其底座同构于 \(A_{Q}\)-模 \(k(s)\)。再证明 \(\mathrm{P}(s)/\mathrm{P}(s)\mathrm{J}_{\mathrm{Q}}\) 同构于 \(k(s)\)。
\item
  假设 Q 中每个圈的长为零。证明对每个 \(s \in \operatorname{Som}(Q)\)，从 \(k\) 到 \(\operatorname{End}(\mathrm{P}(s))\) 的典范同态是一个同构。证明对每个投射且不可分解的 \(A_Q\)-模 P，存在 Q 的唯一顶点 s 使 P 同构于 \(\mathrm{P}(s)\)。
\end{enumerate}

\begin{enumerate}
\def\labelenumi{\arabic{enumi})}
\setcounter{enumi}{7}
\tightlist
\item
  对 \(s \in \operatorname{Som}(Q)\)，令 \(\mathrm{P}(s) = x_s\mathrm{A}_{\mathrm{Q}}\)。设 M 是一个 \(\mathrm{A}_{\mathrm{Q}}\)-模。
\end{enumerate}

\begin{enumerate}
\def\labelenumi{\alph{enumi})}
\tightlist
\item
  设 \(f \in \mathrm{Fl}(\mathbf{Q})\)，令 \(u = o(f)\)，\(v = t(f)\)。证明存在唯一的一对 \((\varphi_f', \varphi_f'')\)（\(\mathrm{A}_{\mathrm{Q}}\)-模同态）
\end{enumerate}

\[
\varphi_ {f} ^ {\prime} \colon \mathrm{M} x _ {u} \otimes_ {k} \mathrm{P} (v) \to \mathrm{M} x _ {u} \otimes_ {k} \mathrm{P} (u), \quad \varphi_ {f} ^ {\prime \prime} \colon \mathrm{M} x _ {u} \otimes_ {k} \mathrm{P} (v) \to \mathrm{M} x _ {v} \otimes_ {k} \mathrm{P} (v),
\]

 使得 \(\varphi_f'(m \otimes p) = m \otimes x_f p\) 且 \(\varphi_f''(m \otimes p) = mx_f \otimes p\)（对一切 \(m \in \mathrm{M}x_u\) 及每个 \(p \in \mathrm{P}(v)\)）。

\begin{enumerate}
\def\labelenumi{\alph{enumi})}
\setcounter{enumi}{1}
\item
  设 \(s \in \text{Som}(Q)\)。证明存在唯一的 \(A_{Q}\)-模同态 \(\psi_{s}: Mx_{s} \otimes_{k} P(s) \to M\)，使得 \(\psi_{s}(m \otimes p) = mp\)（对 \(p \in P(s)\) 与 \(m \in Mx_{s}\)）。
\item
  令 \(\varphi = \bigoplus \varphi_f' - \bigoplus \varphi_f''\)，\(\psi = \sum \psi_s\)。证明有 \(\mathrm{A}_{\mathrm{Q}}\)-模的正合列
\end{enumerate}

\[
0 \to \bigoplus_ {f \in \operatorname{Fl} (\mathrm{Q})} \operatorname{M} x _ {o (f)} \otimes_ {k} \operatorname{P} (t (f)) \xrightarrow {\varphi} \bigoplus_ {s \in \operatorname{Som} (\mathrm{Q})} \operatorname{M} x _ {s} \otimes_ {k} \operatorname{P} (s) \xrightarrow {\psi} \operatorname{M} \to 0.
\]

\(d\)) 证明代数 \(\mathrm{A_Q}\) 的每个右理想都是投射的。

\begin{enumerate}
\def\labelenumi{\arabic{enumi})}
\setcounter{enumi}{8}
\tightlist
\item
  \begin{enumerate}
  \def\labelenumii{\alph{enumii})}
  \tightlist
  \item
    设 M 与 N 是 \(A_Q\)-模；对 \(s \in \operatorname{Som}(Q)\)，令 \(M_s=Mx_s\)，\(N_s=Nx_s\)。对 \(\varphi = (\varphi_s)_{s \in \operatorname{Som}(Q)}\) 与 \(f \in \operatorname{Fl}(Q)\)，如下定义映射 \(\theta_f(\varphi): Mx_{o(f)} \to Nx_{t(f)}\)
  \end{enumerate}
\end{enumerate}

\[
\theta_ {f} (\varphi) (m) = \varphi_ {o (f)} (m) x _ {f} - \varphi_ {t (f)} (m x _ {f}).
\]

证明映射

\[
\theta \colon \prod_ {s \in \operatorname{Som} (\mathrm{Q})} \operatorname{Hom} _ {k} \left(\mathrm{M} _ {s}, \mathrm{N} _ {s}\right)\rightarrow \prod_ {f \in \operatorname{Fl} (\mathrm{Q})} \operatorname{Hom} \left(\mathrm{M} _ {o (f)}, \mathrm{N} _ {t (f)}\right)
\]

 \(\varphi\mapsto(\theta_{f}(\varphi))_{f\in\mathrm{Fl}(Q)}\) 是线性的。

\begin{enumerate}
\def\labelenumi{\alph{enumi})}
\setcounter{enumi}{1}
\item
  证明 \(\operatorname{Ker}(\theta)\) 同构于 \(\operatorname{Hom}_{\mathrm{A}_{\mathrm{Q}}}(\mathrm{M},\mathrm{N})\)，且 \(\operatorname{Coker}(\theta)\) 同构于 \(\operatorname{Ext}_{\mathrm{A}_{\mathrm{Q}}}^{1}(\mathrm{M},\mathrm{N})\)。证明 \(\operatorname{Ext}_{\mathrm{A}_{\mathrm{Q}}}^{p}(\mathrm{M},\mathrm{N}) = 0\)（对每个整数 \(p \geqslant 2\)）。
\item
  设 k 是一个交换域且箭图 Q 是有限的。若 M 与 N 是有限维 k-向量空间，则 \(\operatorname{Ext}_{\mathrm{A}_{\mathrm{Q}}}^{p}(\mathrm{M},\mathrm{N})\) 对每个整数 \(p \geqslant 0\) 是一个有限维 k-向量空间，且当 \(p \geqslant 2\) 时为零，并有
\end{enumerate}

\[
\begin{array}{l} \dim_ {k} (\mathrm{Hom} _ {\mathrm{A} _ {\mathrm{Q}}} (\mathrm{M}, \mathrm{N})) - \dim_ {k} (\mathrm{Ext} _ {\mathrm{A} _ {\mathrm{Q}}} ^ {1} (\mathrm{M}, \mathrm{N})) \\ = \sum_ {s \in \mathrm{Som} (\mathrm{Q})} \dim (\mathrm{M} _ {s}) \dim (\mathrm{N} _ {s}) - \sum_ {f \in \mathrm{Fl} (\mathrm{Q})} \dim (\mathrm{M} _ {o (f)}) \dim (\mathrm{N} _ {t (f)}). \end{array}
\]

\(d\)) 设 \(k\) 是一个交换域。证明对 Q 的每一对顶点 \((u,v)\)，向量空间 \(\mathrm{Ext}_{\mathrm{A_Q}}^1 (k(u),k(v))\) 的维数等于 Q 中以 \(v\) 为起点、\(u\) 为终点的箭的集合的基数。

\begin{enumerate}
\def\labelenumi{\arabic{enumi})}
\setcounter{enumi}{9}
\item
  设 \(k\) 是一个代数闭域。设 Q 与 Q' 是每个圈的长均为零的有限箭图。证明代数 A \(_{Q'}\) 与代数 A \(_{Q}\) Morita 等价（A, VIII, §6）当且仅当箭图 Q 与 Q' 同构。（取 \(k\)-单模 \(k(s)\)（\(s \in \text{Som}(Q)\)）作为模 M、N。）
\end{enumerate}

\exercisegroup

\begin{enumerate}
\def\labelenumi{\arabic{enumi})}
\item
  设 G 是一个图，\(Q = (S, F, o, t)\) 是它的底箭图。
\end{enumerate}

\begin{enumerate}
\def\labelenumi{\alph{enumi})}
\item
  证明有 \(d_{+}(s) = d_{-}(s)\)（对所有 \(s \in S\)）。这时称这个整数为 G 在 s 处的次数，记作 \(d(s)\)。
\item
  设 G 是有限的。为使 G 中存在一条闭道路 \((a_{0}, f_{1}, \ldots, f_{n}, a_{n})\)，使得对 G 的每条边 \(\varphi\)，存在唯一的整数 n 使 \(\varphi = \{f_{n}, \overline{f_{n}}\}\)，必要且充分的是 G 连通且 \(d(s)\) 对所有 \(s \in \text{Som}(G)\) 为偶数。
\end{enumerate}

\begin{enumerate}
\def\labelenumi{\arabic{enumi})}
\setcounter{enumi}{1}
\tightlist
\item
  设 \(\mathrm{G} = (\mathrm{S}, \mathrm{F}, o, t)\) 是一个箭图。若存在 S 的一个分成子集的分割 \(P = \{S_{1}, S_{2}\}\)，使得对每条箭 \(f \in F\)，起点 \(o(f)\) 与终点 \(t(f)\) 位于分割 P 的不同部分，则称 G 是一个二部箭图。
\end{enumerate}

\begin{enumerate}
\def\labelenumi{\alph{enumi})}
\item
  设 G 是连通的。证明如上的分割 P 至多存在一个。
\item
  证明图 G 是二部的当且仅当 G 中每条回路的长均为偶数。
\end{enumerate}

\begin{enumerate}
\def\labelenumi{\arabic{enumi})}
\setcounter{enumi}{2}
\item
  设 G 是一个顶点集为无限集的连通局部有限图。证明存在 G 的箭的一个可合成序列 \((f_{n})_{n\in\mathbf{N}}\)，其诸起点两两不同（König 定理）。
\item
  设 G 是一个有限图，\(Q = (S, F, o, t)\) 是它的底箭图。我们给空间 \(C^{S}\) 赋以唯一的埃尔米特空间结构，使其典范基标准正交。设 L: \(C^{S} \to C^{S}\) 是由 \(\mathrm{L}(u)(x) = \sum_{o(f)=x} (u(t(f)) - u(x))\)（\(x \in S\)，\(u \in C^{S}\)）给出的映射。
\end{enumerate}

\begin{enumerate}
\def\labelenumi{\alph{enumi})}
\item
  证明 \(C^{S}\) 的自同态 L 是正的（EVT, V, p. 45, 定义 6）。
\item
  证明 L 的核由在 G 的每个连通分支上为常值的函数 \(u \in C^{S}\) 组成。
\item
  设 \(F_{1}\) 是 G 的一个定向。设 \(E = (e_{s,f})\) 是如下给出的类型 \((\mathrm{S}, \mathrm{F}_{1})\) 的矩阵
\end{enumerate}

\[
 e _ {s, f} = \left\{ \begin{array}{l l} 1 & \text {if } o (f) = s \text {and } t (f) \neq s; \\ - 1 & \text {if } o (f) \neq s \text {and } t (f) = s; \\ 0 & \text {otherwise. } \end{array} \right.
\]

证明矩阵 \(\Lambda\)（即 L 在 \(\mathbf{C}^{\mathrm{S}}\) 的典范基下的矩阵）由 \(\Lambda = E\cdot {}^{t}E\) 给出。

\(d)\) 设 G 连通且其顶点集非空；设 \(x\) 是 G 的一个顶点，令 \(\mathrm{S}' = \mathrm{S} - \{x\}\)。设 T 是 F 的一个基数为 \(\mathrm{Card}(\mathrm{S}')\) 的子集。证明 T 是 G（赋以定向 \(\mathrm{F}_1\)）的某个极大定向子树的箭集当且仅当矩阵 \(E_{\mathrm{T}}'\)（类型为 \((\mathrm{S}',\mathrm{T})\)，由 \(E\) 导出）的秩等于 \(\mathrm{Card}(\mathrm{S}')\)。由此推出矩阵 \(A_{\mathrm{S}',\mathrm{S}'}\) 的行列式等于 G 的极大定向子树的个数。（利用 A, III, §8, p. 192 习题 6。）

\begin{enumerate}
\def\labelenumi{\alph{enumi})}
\setcounter{enumi}{4}
\tightlist
\item
  设 W 是 \(\mathrm{Ker}(\mathrm{L})\) 的正交补。证明 W 是 \(\mathbf{C}^{\mathrm{S}}\) 的一个在 W 下稳定的子空间，且 \(\det (\mathrm{L}|_{\mathrm{W}})\) 等于图 G 的极大森林的个数（Kirchhoff 定理）。
\end{enumerate}

\begin{enumerate}
\def\labelenumi{\arabic{enumi})}
\setcounter{enumi}{4}
\tightlist
\item
  \begin{enumerate}
  \def\labelenumii{\alph{enumii})}
  \tightlist
  \item
    设 n 是一个自然整数 \(\geqslant 1\)，G 是一个图，其顶点集为集合 \(\mathbf{Z}/n\mathbf{Z}\)，且两个顶点 i 与 j 由一条箭相连当且仅当 \(i-j=\pm1\pmod n\)。明确描述 G 的极大森林的集合。
  \end{enumerate}
\end{enumerate}

\begin{enumerate}
\def\labelenumi{\alph{enumi})}
\setcounter{enumi}{1}
\item
  设 G 是一个图，S 是它的顶点集。假设对 G 的每一对顶点 \((x,y)\)，G 中以 x 为起点、y 为终点的箭的集合在 \(x \neq y\) 时基数为 1，否则为 0（“完全图”）。计算 G 的极大森林的集合的基数。
\item
  设 G 是一个图，S 是它的顶点集。设 \((\mathrm{S}_{1},\mathrm{S}_{2})\) 是 S 的一个分割；假设对 G 的每一对顶点 \((x,y)\)，G 中以 x 为起点、y 为终点的箭的集合在 \((x,y)\in\mathrm{S}_{1}\times\mathrm{S}_{2}\) 时基数为 1，否则为 0（“完全二部图”）。计算 G 的极大森林的集合的基数。
\end{enumerate}

\begin{enumerate}
\def\labelenumi{\arabic{enumi})}
\setcounter{enumi}{5}
\tightlist
\item
  设 \(c\) 是一个自然整数。
\end{enumerate}

\begin{enumerate}
\def\labelenumi{\alph{enumi})}
\tightlist
\item
  存在唯一的映射族 \(\mathrm{R}_{c}\colon(\mathbf{N}^{*})^{c}\to\mathbf{N}^{*}\)（\(c\geqslant2\)），满足下列性质：
\end{enumerate}

\begin{enumerate}
\def\labelenumi{(\roman{enumi})}
\item
  若存在 i 使 \(n_{i}=1\)，则 \(\mathrm{R}_{c}(n_{1},\ldots,n_{c})=1\)；
\item
  我们有 \(\mathrm{R}_{2}(m,n)=\mathrm{R}_{2}(m-1,n)+\mathrm{R}_{2}(m,n-1)\)（若 m 与 n 是 \(\geqslant2\) 的整数）；
\item
  若 \(c \geqslant 3\)，则有 \(\mathrm{R}_c(n_1, \ldots, n_c) = \mathrm{R}_{c-1}(n_1, \ldots, n_{c-2}, \mathrm{R}_2(n_{c-1}, n_c))\)（对所有 \((n_1, \ldots, n_c) \in \mathbf{N}^c\)）。
\end{enumerate}

\begin{enumerate}
\def\labelenumi{\alph{enumi})}
\setcounter{enumi}{1}
\item
  我们有 \(\mathrm{R}_2(m,n) = \binom{m + n - 2}{m - 1}\)（对每对整数 \((m,n)\)，均 \(\geqslant 1\)）。
\item
  对每个有限序列 \((n_{1},\ldots,n_{c})\)（各项为 \(\geqslant1\) 的自然整数）及每个自然整数 \(R\geqslant1\)，我们称性质 \(\boldsymbol{R}(n_{1},\ldots,n_{c};\mathrm{R})\) 成立，如果对每个其顶点集 S 的基数 \(\geqslant R\) 的完全图及其箭集 F 的每个分割 \(P=(F_{1},\ldots,F_{c})\)，存在整数 \(i\in\{1,\ldots,c\}\) 及 S 的基数为 \(n_{i}\) 的子集 A，使得连接 A 中两个元素的 G 的每条箭都属于 \(F_{i}\)。
\end{enumerate}

证明性质 \(\boldsymbol{R}(n_{1},\ldots,n_{c};\mathrm{R}_{c}(n_{1},\ldots,n_{c}))\)。

\begin{enumerate}
\def\labelenumi{\alph{enumi})}
\setcounter{enumi}{3}
\item
  用 \(\mathrm{R}(n_{1},\ldots,n_{c})\) 表示最小的整数 \(R\geqslant1\)，使得性质 \(\boldsymbol{R}(n_{1},\ldots,n_{c};\mathrm{R})\) 成立（“Ramsey 数”）。
\item
  我们有 \(R(n,2)=n\)（对每个整数 \(n\geqslant2\)）；又有 \(R(3,3)=6\)。
\item
  证明 R(4,3) = 9。
\end{enumerate}

计算 R(5, 5)；计算 R(6, 6)。

\begin{enumerate}
\def\labelenumi{\arabic{enumi})}
\setcounter{enumi}{6}
\tightlist
\item
  设 G 是一个右作用于集合 X 上的群，A 是 G 的一个子集。Schreier 图 \(\mathcal{G}_{\mathrm{G}}(\mathrm{X},\mathrm{A})\) 是由以下方式定义的箭图 \((\mathrm{S},\mathrm{F},o,t)\)：
\end{enumerate}

- 其顶点集为 S = X；

- 其箭集为 F = X × A；

- 映射 \(o\) 与 \(t\) 由 \(o(x, a) = x\) 与 \(t(x, a) = x \cdot a\)（对一切 \((x, a) \in \mathrm{X} \times \mathrm{A}\)）给出。

  当 X = G 且 G 通过右乘作用于其上时，我们把这个箭图记作 \(\mathcal{G}(G,A)\)，并称之为 G 相对于 A 的 Cayley 箭图。

我们把与这个箭图相伴的图称为 Schreier 图（相应地，Cayley 图）。

\begin{enumerate}
\def\labelenumi{\alph{enumi})}
\item
  证明箭图 \(\mathcal{G}_{\mathrm{G}}(\mathrm{X},\mathrm{A})\) 连通当且仅当由 A 生成的 G 的子群 H 在 X 上至多有一个轨道。
\item
  设 A 生成 G。证明 Cayley 图 \(\mathcal{G}(\mathrm{G},\mathrm{A})\) 连通，且它是二部的当且仅当存在群同态 \(\varphi \colon \mathrm{G}\to \mathbf{Z} / 2\mathbf{Z}\) 使得 \(\varphi (a) = 1\)（对所有 \(a\in \mathrm{A}\)）。
\item
  证明存在 G 在 Cayley 图 \(\mathcal{G}(\mathrm{G},\mathrm{A})\) 上的唯一作用，使得 G 在该图顶点上诱导的作用是 G 通过左乘作用于其自身的作用。
\end{enumerate}

\begin{enumerate}
\def\labelenumi{\arabic{enumi})}
\setcounter{enumi}{7}
\tightlist
\item
  图 G 中的一条回路 \((a_0, f_1, \ldots, f_n, a_n)\) 称为哈密顿回路，如果对 G 的每个顶点 \(s\)，存在唯一的整数 \(m \in \{1, \ldots, n\}\) 使 \(a_m = s\)。若一个图在 G 中有一条哈密顿回路，则称该图为哈密顿图。
\end{enumerate}

\begin{enumerate}
\def\labelenumi{\alph{enumi})}
\item
  设 S 是一个基数 \(\geqslant 3\) 的有限集，\(K_{S}\) 是以 S 为顶点集的一个完全图。证明图 \(K_{S}\) 是哈密顿图。更精确地说，证明对 G 的每个定向 A，\(K_{S}\) 中存在一条其每条箭都属于 A 的哈密顿回路。
\item
  用 \(G_{S}\) 表示 \(K_{S}\) 的顶点集等于 S 的子图组成的集合。设 \(\preccurlyeq\) 是集合 \(G_{S}\) 上最粗的序关系，使得 \(G \preccurlyeq G'\)（只要存在 S 的元素 u、v，满足 \(\deg(u) + \deg(v) \geqslant \text{Card}(S)\)，且 \(\text{Fl}(G')\) 是 \(\text{Fl}(G)\) 与 \(K_{S}\) 中以 u、v 为端点的两条箭的并）。
\end{enumerate}

  证明对每个 \(G \in G_{S}\)，\(G_{S}\) 中满足 \(G \preccurlyeq H\) 的元素 H 组成的集合有唯一的极大元；把它记作 \(c(G)\)。

\begin{enumerate}
\def\labelenumi{\alph{enumi})}
\setcounter{enumi}{2}
\item
  设 G 是 \(K_{S}\) 的一个以 S 为顶点集的子图。证明图 G 是哈密顿图当且仅当图 \(c(\mathrm{G})\) 也是哈密顿图。
\item
  设 G 是 \(K_{S}\) 的一个子图，其顶点集为 S 且在每个顶点处的次数 \(\geqslant\) Card(S)/2。证明 \(c(\mathrm{G})\) 等于 \(K_{S}\)。由此推出 G 是一个哈密顿图（G. Dirac 定理）。
\end{enumerate}

\begin{enumerate}
\def\labelenumi{\arabic{enumi})}
\setcounter{enumi}{8}
\tightlist
\item
  设 (W, S) 是一个 Coxeter 系（LIE, IV, §1, p. 11, 定义 3）；假设群 W 是有限的。
\end{enumerate}

\begin{enumerate}
\def\labelenumi{\alph{enumi})}
\item
  设 (W, S) 是不可约的且至多有两个顶点。证明 Cayley 图 \(\mathcal{G}(\mathrm{W},\mathrm{S})\) 是哈密顿图。
\item
  设 (W, S) 是不可约的且至少有三个顶点。设 \(s \in S\) 是与 (W, S) 相伴的 Coxeter 图的一个末端顶点（LIE IV, 附录）。令 \(S' = S - \{s\}\)，并用 \(W'\) 表示由 \(S'\) 生成的 W 的子群。设 \(\Gamma\) 是这样的图：其顶点集是 W 模 \(W'\) 的左陪集的集合，且两个不同的顶点 A 与 B 由一条边相连恰好当且仅当 \(\cap\) 非空。利用图 \(\Gamma\) 的一个极大子树，证明若 Cayley 图 \(\mathcal{G}(W', S')\) 是哈密顿图，则图 \(\mathcal{G}(W, S)\) 也是哈密顿图。
\item
  证明 Cayley 图 \(\mathcal{G}(\mathrm{W},\mathrm{S})\) 是哈密顿图（Conway-Sloane-Wilks 定理）。
\end{enumerate}

\begin{enumerate}
\def\labelenumi{\arabic{enumi})}
\setcounter{enumi}{9}
\tightlist
\item
  设 G 是一个连通图，其顶点集为 S。对 \(x, y \in S\)，令
\end{enumerate}

\(d_{\mathrm{G}}(x,y) = \inf \{\ell \in \mathbf{N}\colon \text{存在一条长度为} \ell \text{ 的路径连接 } x \text{ 与 } y\}.\)

\begin{enumerate}
\def\labelenumi{\alph{enumi})}
\item
  证明 \(d_{\mathrm{G}}\) 是集合 S 上的一个距离。
\item
  设 \(G_{1}\) 与 \(G_{2}\) 是图，\(\varphi\colon G_{1}\to G_{2}\) 是一个图态射。证明有 \(d_{\mathrm{G}_{2}}(\varphi(x),\varphi(y))\leqslant d_{\mathrm{G}_{1}}(x,y)\)（对每一对顶点 \((x,y)\)，均属于 \(G_{1}\)）。
\end{enumerate}

\begin{enumerate}
\def\labelenumi{\arabic{enumi})}
\setcounter{enumi}{10}
\tightlist
\item
  设 G 是一个连通图，其顶点集 S 非空。我们称 \(\operatorname{diam}(G)\)（度量空间 (S, \(d_{\mathrm{G}}\)) 的直径）为 G 的直径。
\end{enumerate}

\begin{enumerate}
\def\labelenumi{\alph{enumi})}
\item
  证明 \(\text{diam}(G) \leqslant \text{Card}(S)-1\)。对哪些图该不等式取等号？
\item
  设 v 是一个整数 \(\geqslant 2\)，使得 G 的每个顶点的次数至多为 v。证明 \(\text{Card}(S) \leqslant v^{\text{diam}(G)}\)。对哪些图该不等式取等号？
\item
  设 n 是一个满足 \(n \geqslant 2\) 的整数。设 \(T_{n}\) 是 \(S_{n}\) 中由支集为 \(\{m, m + 1\}\)（\(m \in \{1, \ldots, n - 1\}\)）的对换组成的子集。证明 Cayley 图 \(\mathcal{G}(\mathfrak{S}_{n}, \mathrm{T}_{n})\) 的直径等于 \(n(n - 1)/2\)。
\item
  设 n 是一个满足 \(n \geqslant 2\) 的整数。设 \(\tau\) 是支集为 \(\{1,2\}\) 的对换，\(\sigma\) 是循环 \((1,2,\ldots,n) \mapsto (2,3,\ldots,n,1)\)，它属于群 \(S_n\)。设 \(G_n\) 是 Cayley 图 \(\mathcal{G}(\mathfrak{S}_n,\{\sigma,\tau\})\)。证明 \(n(n-1)/6 \leqslant \text{diam}(G_n) \leqslant 3n(n-1)/2\)。（设 T 是由序列 \((x,y,z)\)（元素取自 \(\{1,\ldots,n\}\)，满足 \(x \textless{} y \textless{} z\)）构成的集合。若 \(t=(x,y,z)\in\text{T}\)，称置换 \(\lambda\in S_n\) 保持 t 的循环次序，如果有 \(\lambda(x)<\lambda(y)<\lambda(z)\)，或 \(\lambda(y)<\lambda(z)<\lambda(x)\)，或 \(\lambda(z)<\lambda(x)<\lambda(y)\)。注意置换 \(\sigma\) 保持 T 的每个元素的循环次序，而置换 \(\tau\) 保持 T 的每个不形如 \((1,2,x)\) 的元素的循环次序。由此推出从恒同置换到置换 \((1,\ldots,n) \mapsto (n,n-1,\ldots,1)\) 的距离至少等于 \(n(n-1)/6\)。）
\end{enumerate}

\begin{enumerate}
\def\labelenumi{\arabic{enumi})}
\setcounter{enumi}{11}
\tightlist
\item
  设 \(n\) 是一个整数 \(\geqslant 2\)，\(p\) 是一个素数 \(>n/2\)。设 H 是 \(\mathfrak{S}_n\) 的一个子群，它传递地作用于 \(\{1, \ldots, n\}\) 上，且包含一个对换和一个 \(p\) 阶循环。设 G 是一个图，满足
\end{enumerate}

- 其顶点集为 \(S = \{1, \ldots, n\}\)；

- 其箭集 F 是由对 \((i,j)\in\mathrm{S}\times\mathrm{S}\)（使对换 \(\tau_{i,j}\) 属于 H）组成的集合；

- 我们有 \(o((i,j)) = i\)，\(t((i,j)) = j\)，且 \(\overline{(i,j)} = (j,i)\)（对一切 \((i,j) \in \mathrm{F}\)）。

\begin{enumerate}
\def\labelenumi{\alph{enumi})}
\item
  证明 G 的每个连通分支都是完全图。
\item
  证明若 G 连通，则 \(H = S_{n}\)。
\item
  证明存在一个从 \(\mathfrak{S}_n\) 到 \(\operatorname{Aut}(\mathbf{G})\) 的群同态，使得对每个 \(\sigma \in \mathfrak{S}_n\)，映射 \(\operatorname{Som}(\sigma)\) 等于置换 \(\sigma\)。然后证明 \(\mathfrak{S}_n\) 传递地作用于 \(\mathbf{G}\) 的连通分支所成的集合上，且这些连通分支两两同构。
\end{enumerate}

\(d)\) 设 \(\sigma\in\mathrm{H}\) 是一个 \(p\) 阶循环。证明 \(\sigma\) 稳定 G 的每个连通分支。由此得出 \(\mathrm{H}=\mathfrak{S}_{n}\)。

\exercisegroup

\begin{enumerate}
\def\labelenumi{\arabic{enumi})}
\tightlist
\item
  设 \((\mathrm{C}_{i})_{i\in\mathrm{I}}\) 是一族范畴。设 C 是诸范畴 \(C_{i}\) 的底箭图组成的族的和箭图。
\end{enumerate}

\begin{enumerate}
\def\labelenumi{\alph{enumi})}
\item
  证明箭图 C 赋以由诸范畴 \(C_{i}\) 的合成法则诱导的合成法则后是一个范畴。我们称之为族 \(C_{i}\) 的和范畴。
\item
  证明从 \(C_{i}\) 到 C 的典范态射是一个函子。
\item
  若诸 \(C_{i}\) 都是群胚，证明范畴 C 是一个群胚。
\end{enumerate}

\begin{enumerate}
\def\labelenumi{\arabic{enumi})}
\setcounter{enumi}{1}
\tightlist
\item
  设 \((\mathrm{C}_{i})_{i\in\mathrm{I}}\) 是一族范畴。设 C 是诸范畴 \(C_{i}\) 的底箭图的积。
\end{enumerate}

\begin{enumerate}
\def\labelenumi{\alph{enumi})}
\item
  证明箭图 C 赋以由诸范畴 \(C_{i}\) 的合成法则诱导的合成法则后是一个范畴。称之为族 \((\mathrm{C}_{i})\) 的积范畴。
\item
  证明典范态射 \(pr_{i}: C \rightarrow C_{i}\) 是一个函子。
\item
  若诸 \(C_{i}\) 都是群胚，证明范畴 C 是一个群胚。
\end{enumerate}

\begin{enumerate}
\def\labelenumi{\arabic{enumi})}
\setcounter{enumi}{2}
\tightlist
\item
  设 C 是一个范畴；用 m 表示它的合成法则。设 C° 是反有向图，赋以由 \(m^{\circ}(f,g)=m(g,f)\)（对 C 的每对可合成的箭 \((f,g)\)）给出的合成法则。
\end{enumerate}

\begin{enumerate}
\def\labelenumi{\alph{enumi})}
\item
  证明 C° 是一个范畴。称之为 C 的反范畴。
\item
  若 C 是一个群胚，证明 C° 是一个同构于 C 的群胚。
\end{enumerate}

\begin{enumerate}
\def\labelenumi{\arabic{enumi})}
\setcounter{enumi}{3}
\tightlist
\item
  我们把群胚 G 的质量 \(\mu(\mathrm{G})\) 定义为诸和 \(\sum_{a\in\mathrm{A}}\mathrm{Card}(\mathrm{G}_{a})^{-1}\) 的上确界，其中 A 取遍 \(\mathrm{Som}(\mathrm{G})\) 中与 G 的每个轨道至多交于一点的子集所成的集合。若 \(\mu(\mathrm{G})\in\mathbf{R}\)，称 G 具有有限质量；若 G 的轨道的集合有限且它在每点的迷向群有限，称 G 是本质有限的。
\end{enumerate}

\begin{enumerate}
\def\labelenumi{\alph{enumi})}
\item
  证明一个本质有限的群胚具有有限质量。
\item
  设 G 是与有限群 \(\Gamma\) 相伴的群胚。它是本质有限的，且有 \(\mu(\mathrm{G}) = 1/\mathrm{Card}(\Gamma)\)。
\item
  设 \(\Gamma\) 是一个有限群，X 是一个赋有 \(\Gamma\) 的右作用的有限集，G 是相伴的群胚。它是本质有限的，且有 \(\mu(\mathrm{G}) = \mathrm{Card}(\mathrm{X}) / \mathrm{Card}(\Gamma)\)。
\end{enumerate}

\(d)\) 设 \((\mathrm{G}_i)_{i\in \mathrm{I}}\) 是一族具有有限质量的群胚。族 \((\mathrm{G}_i)\) 的和群胚与积群胚的质量分别为 \(\sum \mu (\mathrm{G}_i)\) 与 \(\prod \mu (\mathrm{G}_i)\)。

\begin{enumerate}
\def\labelenumi{\alph{enumi})}
\setcounter{enumi}{4}
\tightlist
\item
  设 X 是一个集合。我们定义一个群胚 G，其顶点集是 X 的有限子集组成的集合，且对 X 的每对有限子集 \((Y_{1}, Y_{2})\)，以 \(Y_{1}\) 为起点、\(Y_{2}\) 为终点的箭的集合是从 \(Y_{1}\) 到 \(Y_{2}\) 的双射组成的集合，箭的合成由映射的合成给出。证明 G 是本质有限的，并计算它的质量。
\end{enumerate}

\begin{enumerate}
\def\labelenumi{\arabic{enumi})}
\setcounter{enumi}{4}
\tightlist
\item
  \begin{enumerate}
  \def\labelenumii{\alph{enumii})}
  \tightlist
  \item
    设 \(\varphi\colon G\to G'\) 是一个群胚态射，使得映射 \(\operatorname{Som}(\varphi)\) 是单射。证明 \(\varphi(G)\) 是 \(G'\) 的一个子群胚。
  \end{enumerate}
\end{enumerate}

\begin{enumerate}
\def\labelenumi{\alph{enumi})}
\setcounter{enumi}{1}
\tightlist
\item
  设 G 是一个非传递群胚，其顶点集 \(\operatorname{Som}(G)\) 是二元集 \(\{a,b\}\)。设 G' 是与自由积 \(G_a*G_b\) 相伴的群胚。存在唯一的群胚态射 \(\varphi:G\to G'\)，使得映射 \(\varphi_a\) 与 \(\varphi_b\) 分别是从 \(G_a\) 与 \(G_b\) 到 \(G_a*G_b\) 的典范映射。若 \(G_a\) 或 \(G_b\) 不是平凡群，则该群胚态射的像 \(\varphi(G)\) 不是 \(G'\) 的子群胚。
\end{enumerate}

\begin{enumerate}
\def\labelenumi{\arabic{enumi})}
\setcounter{enumi}{5}
\tightlist
\item
  设 G 与 G' 是图，\(\varphi\colon\mathrm{G}'\to\mathrm{G}\) 是一个图态射。我们称 \((\mathrm{G}',\varphi)\) 是 G 的一个覆叠，如果对每个顶点 \(a\in\mathrm{Som}(\mathrm{G}')\) 及每条箭 \(f\in\mathrm{Fl}(\mathrm{G})\)（以 \(\mathrm{Som}(\varphi)(a)\) 为起点），存在唯一的箭 \(f'\in\mathrm{Fl}(\mathrm{G}')\)（以 \(a\) 为起点）使 \(\mathrm{Fl}(\varphi)(f')=f\)。
\end{enumerate}

我们假设 \((\mathrm{G}^{\prime},\varphi)\) 是 G 的一个覆叠。

\begin{enumerate}
\def\labelenumi{\alph{enumi})}
\item
  证明存在 \(\operatorname{Grp}(G)\) 在 \(\operatorname{Som}(G')\) 上的、相对于映射 \(\operatorname{Som}(\varphi)\) 的唯一作用，使得 \(x \cdot \operatorname{Fl}(\varphi)(f) = t(f)\)（对每个顶点 \(x \in \operatorname{Som}(G')\) 及每条以 x 为起点的箭 \(f \in \operatorname{Fl}(G')\)）。
\item
  由此推出，若 x 与 y 是 G 的属于同一连通分支的顶点，则 \(\text{Card}(\text{Som}(\varphi)^{-1}(x)) = \text{Card}(\text{Som}(\varphi)^{-1}(y))\)。
\end{enumerate}

\begin{enumerate}
\def\labelenumi{\arabic{enumi})}
\setcounter{enumi}{6}
\tightlist
\item
  设 G 是一个图，\((G_{1},\varphi_{1})\) 与 \((G_{2},\varphi_{2})\) 是图 G 的覆叠。
\end{enumerate}

\begin{enumerate}
\def\labelenumi{\alph{enumi})}
\item
  设 \(\psi\colon\mathrm{G}_{1}\to\mathrm{G}_{2}\) 是一个满足 \(\varphi_{2}\circ\psi=\varphi_{1}\) 的图态射。证明有 \(\psi(x\cdot\gamma)=\psi(x)\cdot\gamma\)（对所有 \(x\in\mathrm{Som}(\mathrm{G}_{1})\) 及每条箭 \(\gamma\)（属于 \(\varpi_{\mathrm{G}}\)，起点为 \(\varphi_{1}(x)\)））。
\item
  设 \(u \colon \operatorname{Som}(\mathrm{G}_1) \to \operatorname{Som}(\mathrm{G}_2)\) 是一个映射，满足 \(u(x \cdot \gamma) = u(x) \cdot \gamma\)（对所有 \(x \in \operatorname{Som}(\mathrm{G}_1)\) 及每条箭 \(\gamma\)（属于 \(\varpi_{\mathrm{G}}\)，起点为 \(\varphi_1(x)\)））。证明存在唯一的图态射 \(\psi \colon \mathrm{G}_1 \to \mathrm{G}_2\)，使得 \(\varphi_2 \circ \psi = \varphi_1\) 且 \(\operatorname{Som}(\psi) = u\)。
\item
  设 X 是一个集合，\(p: X \to \text{Som}(G)\) 是一个映射。假设给定了 \(\varpi_{G}\) 相对于 p 在 X 上的一个作用。证明存在一个以 X 为顶点集的图 \(G'\) 及一个图态射 \(\varphi: G' \to G\)，使得 \(\text{Som}(\varphi) = p\)，\((G', \varphi)\) 是 G 的一个覆叠，且 \(\varpi_{G}\) 在 \(\text{Som}(G')\) 上的作用与给定的作用一致。
\end{enumerate}

\begin{enumerate}
\def\labelenumi{\arabic{enumi})}
\setcounter{enumi}{7}
\tightlist
\item
  设 G 是一个连通图，\(a\) 是 G 的一个顶点。
\end{enumerate}

\begin{enumerate}
\def\labelenumi{\alph{enumi})}
\tightlist
\item
  设 \((\mathrm{G}_{1},\varphi_{1})\) 与 \((\mathrm{G}_{2},\varphi_{2})\) 是图 G 的覆叠。设 \(u\colon\operatorname{Som}(\mathrm{G}_{1})\to\operatorname{Som}(\mathrm{G}_{2})\) 是一个满足 \(u(x)\cdot\gamma=u(x\cdot\gamma)\) 的映射
\end{enumerate}

（对所有 \(\gamma \in \varpi_{\mathrm{G},a}\) 及每个顶点 \(x\)（属于 \(\mathrm{G}_1\) 且满足 \(\operatorname{Som}(\varphi_1)(x) = a\)））。证明存在唯一的图态射 \(\psi: \mathrm{G}_1 \to \mathrm{G}_2\) 使得 \(\psi = \operatorname{Som}(u)\)。

\begin{enumerate}
\def\labelenumi{\alph{enumi})}
\setcounter{enumi}{1}
\tightlist
\item
  设 X 是一个赋有群 \(\varpi_{\mathrm{G},a}\) 的作用的集合。证明存在图 G 的一个覆叠 \((\mathrm{G}',\varphi)\)，使得 \(\operatorname{Som}(\varphi)^{-1}(a) = \mathrm{X}\)，且 \(\varpi_{\mathrm{G}}\) 在 \(\operatorname{Som}(\mathrm{G}')\) 上的作用诱导出 \(\varpi_{\mathrm{G},a}\) 在 X 上给定的作用。
\end{enumerate}

\begin{enumerate}
\def\labelenumi{\arabic{enumi})}
\setcounter{enumi}{8}
\tightlist
\item
  设 G 与 G' 是图，\(\varphi\colon G'\to G\) 是一个图态射，使得 \((G',\varphi)\) 是 G 的一个覆叠。
\end{enumerate}

\begin{enumerate}
\def\labelenumi{\alph{enumi})}
\item
  证明图 \(\varphi(\mathrm{Grp}(\mathrm{G}^{\prime}))\) 是群胚 \(\mathrm{Grp}(\mathrm{G})\) 的一个子群胚。
\item
  证明对每个顶点 \(a\)（属于 \(\mathbf{G}'\)），群同态 \(\pi_1(\varphi, a)\) 是单射。
\item
  设 G 连通，并固定 G 的一个顶点 a。设 K 是 \(\varpi_{G,a}\) 的一个子群。证明存在图 G 的一个覆叠 \((H,\psi)\) 及 H 的一个顶点 b，使得 \(\psi(b)=a\) 且 K 是态射 \(\pi_{1}(\psi,a)\) 的像。
\item
  证明自由群的子群是自由的（另见 IV, p. 417, 例 2）。
\item
  设 n 与 q 是整数；若 F 是 n 个生成元上的自由群，而 K 是 F 的指数为 q 的子群，证明 K 是 \(1 + q(n - 1)\) 个生成元上的自由群。
\end{enumerate}

\begin{enumerate}
\def\labelenumi{\arabic{enumi})}
\setcounter{enumi}{9}
\tightlist
\item
  设 F 是一个集合，C 是 \(F \times F\) 的一个子集，\(m: C \to F\) 是一个映射。我们作如下假设：
\end{enumerate}

\begin{enumerate}
\def\labelenumi{(\roman{enumi})}
\item
  由 \((f,g)\mapsto(f,m(f,g))\) 与 \((f,g)\mapsto(m(f,g),g)\) 给出的从 C 到 F × F 的映射都是单射；
\item
  若 \(f, g, h \in \mathrm{F}\) 使得对 \((f, g)\) 与 \((g, h)\) 属于 C，则对 \((f, m(g, h))\) 与 \((m(f, g), h)\) 属于 C，且有 \(m(f, m(g, h)) = m(m(f, g), h)\)。
\item
  对每个 \(f \in \mathrm{F}\)，存在元素 \(o(f)\)、\(t(f)\) 与 \(\overline{f} \in \mathrm{F}\)（由 (i) 知必然唯一），使得对 \((f, t(f))\)、\((o(f), f)\)、\((f, \overline{f})\) 与 \((\overline{f}, f)\) 属于 C，且有 \(m(f, t(f)) = f = m(o(f), f)\) 及 \(m(f, \overline{f}) = o(f)\)。
\end{enumerate}

\begin{enumerate}
\def\labelenumi{\alph{enumi})}
\item
  证明对每个 \(f \in F\)，有 \((\overline{f}, f) \in C\) 且 \(m(\overline{f}, f) = t(f)\)。由此推出 \((o(f), o(f)) \in C\) 且 \(m(o(f), o(f)) = o(f)\)。再证明 \(o(\overline{f}) = t(f)\) 且 \(\overline{\overline{f}} = f\)。
\item
  设 S 是 F 中元素 \(e\)（满足 \((e,e)\in\mathrm{C}\) 且 \(m(e,e)=e\)）组成的集合；用 \(o\) 与 \(t\) 表示由 \(f\mapsto o(f)\) 与 \(f\mapsto t(f)\) 给出的从 F 到 S 的映射。
\end{enumerate}

证明四元组 \(\Gamma = (S, F, o, t)\) 是一个箭图，且映射 \(m\) 是 \(\Gamma\) 上的一个使它成为群胚的合成法则。（这一构造属于 H. BRANDT；参见 “Über eine Verallgemeinerung des Gruppenbegriffes”, Mathematische Annalen (1926), vol. 96, p. 360–366。）

\begin{enumerate}
\def\labelenumi{\arabic{enumi})}
\setcounter{enumi}{10}
\tightlist
\item
  设 G 是一个群胚；设 C 是 G 的可合成的箭对组成的集合，\(m\colon\mathrm{C}\to\mathrm{Fl}(\mathrm{G})\) 是 G 的合成法则。
\end{enumerate}

\begin{enumerate}
\def\labelenumi{\alph{enumi})}
\item
  证明 C 与 m 满足习题 10 的假设 (i)、(ii)、(iii)。用 Γ 表示该习题中的构造所给出的群胚。
\item
  证明由映射 \(e \mapsto o(e)\)（从 \(\operatorname{Som}(\Gamma)\) 到 \(\operatorname{Som}(\mathrm{G})\)）与集合 \(\operatorname{Fl}(\mathrm{G})\) 到其自身的恒同映射组成的对是一个从 \(\Gamma\) 到 G 的群胚同构。
\end{enumerate}

这表明一个群胚由它的箭集及其合成法则确定。

\begin{enumerate}
\def\labelenumi{\arabic{enumi})}
\setcounter{enumi}{11}
\tightlist
\item
  设 G 是一个群，S 是 G 的子群组成的集合，F 是 G 中满足存在 G 的子群 H 与 G 的元素 g 使 X = gH 的子集 X 组成的集合。
\end{enumerate}

\begin{enumerate}
\def\labelenumi{\alph{enumi})}
\item
  设 \(\mu\) 是从 \(G \times G \times G\) 到 \(G\) 的、由 \(\mu(x,y,z)=xy^{-1}z\) 定义的映射。证明一个非空子集 \(X\)（\(G\) 的）属于集合 F 当且仅当 \(\mu(X\times X\times X)\subset X\)。
\item
  证明取 X 的任意元素 g 并令 \(o(\mathrm{X}) = g^{-1}\mathrm{X}\) 与 \(t(\mathrm{X}) = \mathrm{X}g^{-1}\)，便定义了映射 \(o: F \to S\) 与 \(t: F \to S\)。
\item
  若 \(X_{1}\) 与 \(X_{2}\) 是 S 中满足 \(t(\mathrm{X}_{1}) = o(\mathrm{X}_{2})\) 的元素，令 \(m(\mathrm{X}_{1}, \mathrm{X}_{2}) = \mathrm{X}_{1}\mathrm{X}_{2}\)。证明映射 m 是箭图 (S,F,o,t) 上的一个使它成为群胚的合成法则（“Baer 群胚”）。把它记作 \(\mathrm{Ba}(\mathrm{G})\)。
\item
  证明 \(\mathrm{Ba}(\mathrm{G})\) 的轨道是 G 的子群的共轭类。
\item
  证明 Ba(G) 在 S 的元素 H 处的迷向群同构于商群 N(H)/H。
\end{enumerate}

\setcounter{exercisegroup}{4}

\exercisegroup

\begin{enumerate}
\def\labelenumi{\arabic{enumi})}
\tightlist
\item
  设 G 是一个群胚，X 是一个集合，f 是从 Som(G) 到 X 的一个映射。
\end{enumerate}

\begin{enumerate}
\def\labelenumi{\alph{enumi})}
\item
  证明存在一个对 \((K,\alpha)\)，其中 K 是以 X 为顶点集的群胚，\(\alpha: G \rightarrow K\) 是满足 \(\mathrm{Som}(\alpha)=f\) 的群胚态射，且满足下列泛性质：对每个以 X 为顶点集的群胚 H 及每个态射 \(\varphi:G\to H\)（满足 \(\mathrm{Som}(\varphi)=f\)），存在唯一的群胚态射 \(\psi:K\to H\) 使 \(\alpha\circ\psi=\varphi\)。（若 S 是一个集合，用 \(S_{d}\) 表示一个以 S 为顶点集、其轨道与迷向群都化为单独一个元素的群胚。设 \(G'\) 是族 \((\mathrm{G},\mathrm{X}_{\mathrm{d}})\) 的不交并群胚，设 \(\varphi\) 与 \(\psi\) 是从群胚 \(\mathrm{Som}(\mathrm{G})_{\mathrm{d}}\) 到 \(G'\) 的态射，分别满足 \(\mathrm{Som}(\varphi)\) 是 \(\mathrm{Som}(\mathrm{G})\) 到 \(\mathrm{Som}(\mathrm{G}')\) 中的典范单射与 \(\mathrm{Som}(\psi)=f\)。令 \(K=\mathrm{Coeg}(\varphi,\psi).\)。）
\item
  由此推出，对每个以 X 为顶点集的群胚 H，有从 \(\operatorname{Hom}(G, f^{*}H)\) 到 \(\operatorname{Hom}(K, H)\) 上的一个双射。
\end{enumerate}

\begin{enumerate}
\def\labelenumi{\arabic{enumi})}
\setcounter{enumi}{1}
\tightlist
\item
  设 H 与 G 是群，\(\mathcal{H}\) 与 \(\mathcal{G}\) 是与之相伴的群胚；设 \(\varphi\) 与 \(\psi\) 是从 \(\mathcal{H}\) 到 \(\mathcal{G}\) 的群胚态射。
\end{enumerate}

\begin{enumerate}
\def\labelenumi{\alph{enumi})}
\item
  证明对 \((\varphi,\psi)\) 的上同伦子（相应地，余均衡子）的顶点集化为单独一个元素；用 K（相应地，L）表示它在该点的迷向群。
\item
  证明群 K 是自由积 \(G * \mathbf{Z}\) 对 \(G * \mathbf{Z}\) 中包含元素 \(\varphi(g)t\psi(g)^{-1}t^{-1}\)（\(g \in G\)）的最小正规子群所作的商，其中 t 表示 \(\mathbf{Z}\) 的元素 1。
\item
  证明群 L 是群 G 对包含元素 \(\varphi(g)\psi(g)^{-1}\)（\(g \in G\)）的最小正规子群所作的商。
\end{enumerate}

\begin{enumerate}
\def\labelenumi{\arabic{enumi})}
\setcounter{enumi}{2}
\tightlist
\item
  设 I 是一个预序集，\((\mathrm{G}_{i})_{i\in\mathrm{I}}\) 是一族群胚。对 I 中每对元素 \((i,j)\)（满足 \(i\leqslant j\)），设 \(\varphi_{ji}\) 是从 \(G_{i}\) 到 \(G_{j}\) 的群胚态射。假设诸 \(\varphi_{ji}\) 满足下列条件（参见 E, III, p. 61）：
\end{enumerate}

\begin{enumerate}
\def\labelenumi{(\roman{enumi})}
\item
  关系 \(i \leqslant j \leqslant k\) 蕴含 \(\varphi_{ki} = \varphi_{kj} \circ \varphi_{ji}\)；
\item
  对每个 \(i \in \mathrm{I}\)，\(f_{ii}\) 是 \(\mathbf{G}_i\) 的恒同态射。
\end{enumerate}

设 G 是族 \((\mathrm{G}_{i})\) 的和群胚；设 H 是族 \((\mathrm{H}_{ji})_{i\leqslant j}\) 的和群胚，其中已令 \(H_{ji} = G_{i}\)（对每对 \((i,j)\)）。设 \(\varphi\)（相应地，\(\psi\)）是由族 \((\varphi_{ji})_{i\leqslant j}\)（相应地，\((\mathrm{Id}_{\mathrm{G}_{i}})_{i\leqslant j}\)）诱导的从 H 到 G 的态射，并设 C 是对 \((\varphi,\psi)\) 的余均衡子。对 \(i \in I\)，设 \(\varphi_{i}\) 是从 \(G_{i}\) 到 C 的、由 \(G_{i}\) 到 G 的典范态射与 G 到 C 的典范态射复合而成的态射。

\begin{enumerate}
\def\labelenumi{\alph{enumi})}
\tightlist
\item
  证明有 \(\varphi_j \circ \varphi_{ji} = \varphi_i\)（对每对 \((i,j)\)，满足 \(i \leqslant j\)）。b) 证明对每个群胚 D 及每个族 \((g_i)\)（其中 \(g_i\) 是从 \(\mathrm{G}_i\) 到 D 的态射，使得有 \(g_j \circ \varphi_{ji} = g_i\)（对每对 \((i,j)\)，\(i \leqslant j\)）），存在唯一的群胚态射 \(g: \mathrm{C} \to \mathrm{D}\) 使得 \(g \circ \varphi_i = g_i\)（对所有 \(i \in \mathrm{I}\)）。
\end{enumerate}

\chapter{同伦与庞加莱群胚}

\section{同伦，同伦等价}

在本节及随后各节中，I 表示 R 的区间 {[}0,1{]}。

\subsection{同伦的连续映射}

定义 1 —— 设 X 与 Y 是拓扑空间，f 与 g 是从 X 到 Y 的连续映射。连接 f 与 g 的同伦是一个连续映射 \(\sigma: X \times I \to Y\)，使得对每个 \(x \in X\)，有 \(\sigma(x,0) = f(x)\) 与 \(\sigma(x,1) = g(x)\)。若存在连接 f 与 g 的同伦，则称 f 与 g 同伦。

我们称 \(f\) 为起点，\(g\) 为同伦 \(\sigma\) 的终点（参见下文 III, p. 257, 注 2）。

设 A 是 X 的一个子集。若同伦 \(\sigma\) 满足：对每个 \(a \in \mathrm{A}\)，从 I 到 Y 的映射 \(t \mapsto \sigma(a, t)\) 是常值的，则称它在 A 上是固定的。此时，\(\sigma\) 的起点与终点在 A 的每点处取相同的值。

设 X 与 Y 是拓扑空间。我们称同伦 \(\sigma\colon X\times I\to Y\) 与 \(\tau\colon X\times I\to Y\) 是可拼接的，如果 \(\sigma\) 的终点是 \(\tau\) 的起点，亦即有 \(\sigma(x,1)=\tau(x,0)\)（对所有 \(x\in X\)）。此时，映射 \(\sigma*\tau\)（从 \(X\times I\) 到 Y）由下式定义：

\[
(\sigma * \tau) (x, t) = \left\{ \begin{array}{l l} \sigma (x, 2 t) & \text {for} 0 \leqslant t \leqslant 1 / 2 \\ \tau (x, 2 t - 1) & \text {for} 1 / 2 \leqslant t \leqslant 1 \end{array} \right.\tag{1}
\]

它是连续的（TG, I, p. 19, 命题 4），且是连接 \(\sigma\) 的起点与 \(\tau\) 的终点的一个同伦。称之为同伦 \(\sigma\) 与 \(\tau\) 的拼接同伦。

若 \(\sigma\colon X\times\mathbf{I}\to Y\) 是一个同伦，则映射 \(\overline{\sigma}\colon X\times\mathbf{I}\to Y\)（由 \((x,t)\mapsto\sigma(x,1-t)\) 定义）是连接 \(\sigma\) 的终点与 \(\sigma\) 的起点的同伦。我们有 \(\overline{\overline{\sigma}}=\sigma\)。若 \(\sigma\) 与 \(\tau\) 是从 \(X\times\mathbf{I}\) 到 Y 的可拼接同伦，则同伦 \(\overline{\tau}\) 与 \(\overline{\sigma}\) 可拼接，且有 \(\overline{\sigma*\overline{\tau}}=\overline{\tau}*\overline{\sigma}\)。

命题 1 —— 设 X 与 Y 是拓扑空间。关系“f 与 g 同伦”是从 X 到 Y 的连续映射所成的集合 \(\mathcal{C}(\mathrm{X};\mathrm{Y})\) 上的一个等价关系。

设 \(f\) 是 \(\mathcal{C}(\mathrm{X};\mathrm{Y})\) 的一个元素。映射 \(f\circ\mathrm{pr}_{1}\colon\mathrm{X}\times\mathrm{I}\to\mathrm{Y}\) 是连接 \(f\) 与 \(f\) 的同伦；因此该关系是自反的。

设 \(f\) 与 \(g\) 是 \(\mathcal{C}(\mathrm{X};\mathrm{Y})\) 的元素，\(\sigma\colon\mathrm{X}\times\mathbf{I}\to\mathrm{Y}\) 是连接 \(f\) 与 \(g\) 的同伦。于是映射 \(\overline{\sigma}\colon\mathrm{X}\times\mathbf{I}\to\mathrm{Y}\) 是连接 \(g\) 与 \(f\) 的同伦；因此所考虑的关系是对称的。

最后证明它是传递的。若 \(f\)、\(g\)、\(h\) 是 \(\mathcal{C}(\mathrm{X};\mathrm{Y})\) 的元素，\(\sigma\) 是连接 \(f\) 与 \(g\) 的同伦，\(\tau\) 是连接 \(g\) 与 \(h\) 的同伦，则 \(\sigma\) 与 \(\tau\) 可拼接，且 \(\sigma*\tau\) 是连接 \(f\) 与 \(h\) 的同伦。

设 X 与 Y 是拓扑空间。\(\mathcal{C}(\mathrm{X};\mathrm{Y})\) 中的等价关系“f 与 g 同伦”（命题 1）称为同伦关系。\(\mathcal{C}(\mathrm{X};\mathrm{Y})\) 对该关系的商集记作 {[}X;Y{]}。它的元素称为从 X 到 Y 的连续映射的同伦类。连续映射 \(f\colon \mathrm{X}\to \mathrm{Y}\) 的同伦类常记作 {[}f{]}。

命题 2 —— 设 X、Y、Z 是拓扑空间，f 与 \(f'\) 是从 X 到 Y 的连续映射，g 与 \(g'\) 是从 Y 到 Z 的连续映射。若 f 与 \(f'\) 同伦且 g 与 \(g'\) 同伦，则 \(g \circ f\) 与 \(g' \circ f'\) 同伦。

设 \(\sigma\) 是连接 \(f\) 与 \(f'\) 的同伦，\(\tau\) 是连接 \(g\) 与 \(g'\) 的同伦。于是映射 \(\theta\colon X\times I\to Z\)（由 \(\theta(x,t)=\tau(\sigma(x,t),t)\) 定义）是连接 \(g\circ f\) 与 \(g'\circ f'\) 的同伦。

设 X、Y、Z 是拓扑空间。给定同伦类 \(\varphi \in [\mathrm{X};\mathrm{Y}]\) 与 \(\psi \in [\mathrm{Y};\mathrm{Z}]\)，诸映射 \(g \circ f: \mathrm{X} \to \mathrm{Z}\)（其中 \(f \in \varphi\)，\(g \in \psi\)）都属于同一个同伦类（命题 2），把它记作 \(\psi \circ \varphi\)，并称为类 \(\psi\) 与 \(\varphi\) 的复合同伦类。映射 {[}X;Y{]} \(\times\) {[}Y;Z{]} \(\to\) {[}X;Z{]}（它把 \(\psi \circ \varphi\) 赋予 \((\varphi,\psi)\)）称为合成映射。

设 \(\varphi\in[\mathrm{X};\mathrm{Y}]\)；我们有 \(\varphi=\varphi\circ[\mathrm{Id}_{\mathrm{X}}]=[\mathrm{Id}_{\mathrm{Y}}]\circ\varphi\)。

设 X、Y、Z、T 是拓扑空间，\(\varphi \in [\mathrm{X};\mathrm{Y}]\)、\(\psi \in [\mathrm{Y};\mathrm{Z}]\)、\(\chi \in [\mathrm{Z};\mathrm{T}]\)。同伦类 \(\chi \circ (\psi \circ \varphi)\) 等于 \((\chi \circ \psi) \circ \varphi\)；把它记作 \(\chi \circ \psi \circ \varphi\)。

命题 3 —— 设 X 是拓扑空间，\((\mathrm{Y}_{j})_{j\in\mathrm{J}}\) 是一族拓扑空间。从 \([X;\prod_{j\in\mathrm{J}}Y_{j}]\) 到积集 \(\prod_{j\in\mathrm{J}}[X;Y_{j}]\) 的、由 \(\varphi\mapsto([\mathrm{pr}_{j}]\circ\varphi)_{j\in\mathrm{J}}\) 定义的映射是双射。

满射性由 I, p. 25, 命题 1 立得。我们来证明单射性。设 \(f\) 与 \(g\) 是从 X 到 \(\prod_{j \in \mathrm{J}} \mathrm{Y}_j\) 的连续映射。对所有 \(j \in \mathrm{J}\)，令 \(f_j = \mathrm{pr}_j \circ f\)，\(g_j = \mathrm{pr}_j \circ g\)；设 \(f_j\) 与 \(g_j\) 同伦，\(\sigma_j\) 是连接 \(f_j\) 与 \(g_j\) 的同伦。映射 \(\sigma = (\sigma_j)\)（从 \(\mathrm{X} \times \mathbf{I}\) 到 \(\prod_{j \in \mathrm{J}} \mathrm{Y}_j\)）是连续的（出处同上）；它是连接 \(f\) 与 \(g\) 的同伦，由此得命题。

推论 —— 设 \((\mathrm{X}_{j})_{j\in\mathrm{J}}\) 与 \((\mathrm{Y}_{j})_{j\in\mathrm{J}}\) 是有同一指标集的拓扑空间族。对每个 \(j\in J\)，设 \(f_{j}\) 与 \(g_{j}\) 是从 \(X_{j}\) 到 \(Y_{j}\) 的连续映射。若对每个 \(j\in J\)，映射 \(f_{j}\) 与 \(g_{j}\) 都同伦，则积映射 \(f\colon(x_{j})\mapsto(f_{j}(x_{j}))\) 与 \(g\colon(x_{j})\mapsto(g_{j}(x_{j}))\)（从 \(\prod_{j\in J}X_{j}\) 到 \(\prod_{j\in J}Y_{j}\)）也同伦。

假设有 \([f_j] = [g_j]\)（对所有 \(j \in \mathbf{J}\)）。我们有 \([\mathrm{pr}_j] \circ [f] = [f_j \circ \mathrm{pr}_j]\) 与 \([\mathrm{pr}_j] \circ [g] = [g_j \circ \mathrm{pr}_j]\)；于是由命题 2 得 \([\mathrm{pr}_j] \circ [f] = [\mathrm{pr}_j] \circ [g]\)，再由命题 3 得 \([f] = [g]\)。

\subsection{带基点同伦}

设 X 与 Y 是拓扑空间，x 是 X 的一点。若同伦 \(\sigma\colon X \times I \to Y\) 在 \(\{x\}\) 上是固定的，即若从 I 到 Y 的映射 \(t \mapsto \sigma(x, t)\) 是常值的，则称它在 x 处带基点。两个在 x 处带基点的同伦的拼接在 x 处带基点。在 x 处带基点的同伦的起点与终点在 x 处取同一值 y；因此它们是从 \((X, x)\) 到 \((Y, y)\) 的带基点连续映射。

设 \((X,x)\) 与 \((Y,y)\) 是带基点拓扑空间（I, p. 120, 定义 1）。关系“f 由在 x 处带基点的同伦连接到 g”是集合 \(\mathcal{C}((X,x);(Y,y))\) 上的一个等价关系，称为带基点同伦关系。\(\mathcal{C}((X,x);(Y,y))\) 对该等价关系的商集记作 \([(X,x);(Y,y)]\)。它的元素称为从 \((X,x)\) 到 \((Y,y)\) 的带基点连续映射的带基点同伦类。

\pandocbounded{\includegraphics[keepaspectratio]{images/part002_b3045f29.jpg}}

集合 \(\mathcal{C}((\mathrm{X},x);(\mathrm{Y},y))\) 是 \(\mathcal{C}(\mathrm{X};\mathrm{Y})\) 的一个子集。注意，带基点同伦是比同伦在 \(\mathcal{C}(\mathrm{X};\mathrm{Y})\) 上诱导的关系更细的等价关系。一般而言它是一个严格更细的关系（见 III, p. 321, 习题 1 与 III, p. 234, 例 3）。

再设 \((Z,z)\) 是带基点拓扑空间。设 f 与 \(f'\) 是从 \((X,x)\) 到 \((Y,y)\) 的带基点连续映射，g 与 \(g'\) 是从 \((Y,y)\) 到 \((Z,z)\) 的带基点连续映射。若 f 与 \(f'\) 由在 x 处带基点的同伦 \(\sigma\) 相联系，且 g 与 \(g'\) 由在 y 处带基点的同伦 \(\tau\) 相联系，则 \(g \circ f\) 与 \(g' \circ f'\) 由同伦 \(\theta: X \times I \to Z\)，\((u,t) \mapsto \tau(\sigma(u,t),t)\) 相连接，该同伦在 x 处带基点。与上文一样，由此可以定义从 \([(X,x);(Y,y)] \times [(Y,y);(Z,z)]\) 到 \([(X,x);(Z,z)]\) 的合成映射。我们让读者自己陈述并证明关于带基点映射与带基点同伦的、与命题 3 及其推论相仿的命题。

\subsection{同伦等价的空间}

定义 2 —— 设 X 与 Y 是拓扑空间。同伦类 \(\varphi \in [\mathrm{X};\mathrm{Y}]\) 称为可逆的，如果存在同伦类 \(\psi \in [\mathrm{Y};\mathrm{X}]\)，使得 \(\psi \circ \varphi = [\mathrm{Id}_{\mathrm{X}}]\) 且 \(\varphi \circ \psi = [\mathrm{Id}_{\mathrm{Y}}]\)。若一个连续映射的同伦类可逆，则称该连续映射为同伦等价。

设 \(\varphi\in[X;Y]\) 是一个可逆同伦类。具有定义 2 中诸性质的同伦类 \(\psi\in[Y;X]\) 是唯一的。事实上，设 \(\psi,\psi'\) 是具有这些性质的类；我们有

\[
\psi = \psi \circ [ \mathrm{Id} _ {\mathrm{Y}} ] = \psi \circ \varphi \circ \psi^ {\prime} = [ \mathrm{Id} _ {\mathrm{X}} ] \circ \psi^ {\prime} = \psi^ {\prime}.
\]

这个唯一的类称为同伦类 \(\varphi\) 的逆，记作 \(\varphi^{-1}\)。

设 Z 是拓扑空间。我们有 \(\chi\circ\varphi\circ\varphi^{-1}=\chi\)（对所有 \(\chi\in[Y;Z]\)），且 \(\theta\circ\varphi^{-1}\circ\varphi=\theta\)（对所有 \(\theta\in[X;Z].\)）。由此推出，映射 \(\chi\mapsto\chi\circ\varphi\)（从 {[}Y;Z{]} 到 {[}X;Z{]}）与映射 \(\theta\mapsto\theta\circ\varphi^{-1}\)（从 {[}X;Z{]} 到 {[}Y;Z{]}）互为逆双射。同样，映射 \(\chi\mapsto\varphi\circ\chi\)（从 {[}Z;X{]} 到 {[}Z;Y{]}）是双射，其逆是映射 \(\theta\mapsto\varphi^{-1}\circ\theta\)（从 {[}Z;Y{]} 到 {[}Z;X{]}）。

设 X、Y、Z 是拓扑空间，\(\varphi \in [X;Y]\) 与 \(\psi \in [Y;Z]\) 是可逆同伦类。于是类 \(\psi \circ \varphi\) 可逆，其逆为 \(\varphi^{-1} \circ \psi^{-1}\)。事实上，我们有

\[
\begin{array}{r l} & (\psi \circ \varphi) \circ (\varphi^ {- 1} \circ \psi^ {- 1}) = \psi \circ (\varphi \circ \varphi^ {- 1}) \circ \psi^ {- 1} \\ & \qquad = \psi \circ [ \mathrm{Id} _ {\mathrm{Y}} ] \circ \psi^ {- 1} \\ & \qquad = \psi \circ \psi^ {- 1} = [ \mathrm{Id} _ {\mathrm{Z}} ] \end{array}
\]

类似地，\((\varphi^{-1} \circ \psi^{-1}) \circ (\psi \circ \varphi) = [\mathrm{Id}_{\mathrm{X}}]\)。

设 X 与 Y 是拓扑空间，\(f: X \to Y\) 是一个同伦等价。凡其类为 f 的类之逆的连续映射 \(g: Y \to X\)，均称为 f 的同伦逆（或逆）。这样的映射 g 是一个同伦等价。

同胚 \(f\) 是同伦等价，且 \([f]^{-1} = [f^{-1}]\)。两个同伦等价的复合是同伦等价。关系“X 与 Y 是拓扑空间且存在从 X 到 Y 的同伦等价”是一个等价关系。

定义 3 —— 若存在从 X 到 Y 的同伦等价，则称拓扑空间 X 与 Y 同伦等价。

例 —— 1) 空拓扑空间只与它自身同伦等价。

\begin{enumerate}
\def\labelenumi{\arabic{enumi})}
\setcounter{enumi}{1}
\item
  设 X 是非空拓扑空间。X 与一个化为一点的空间同伦等价的充要条件是：存在与 X 的恒同映射同伦的常值映射 \(p: X \rightarrow X\)。（此时对每个常值映射 \(q \colon X \to X\) 亦然，因为 \([p] = [p] \circ [q] = [\mathrm{Id}_{\mathrm{X}}] \circ [q] = [q]\)。）事实上，设 \(P\) 是化为一点的空间，\(f \colon P \to X\) 是一个映射，\(g \colon X \to P\) 是从 X 到 P 的唯一映射。我们有 \(g \circ f = \mathrm{Id}_{\mathrm{P}}\)；f 是同伦等价的充要条件是 \(f \circ g\) 与 \(\mathrm{Id}_{\mathrm{X}}\) 同伦。而 \(f \circ g\) 是常值的，其像为 \(f(P)\)。
\item
  若存在在 x 处带基点的同伦，把 \(Id_{X}\) 连接到从 X 到 X、以 \(\{x\}\) 为像的常值映射，则称带基点拓扑空间 (X,x) 是可缩的，或称拓扑空间 X 在 x 处可缩。这样的空间与一点同伦等价。然而，存在与一点同伦等价但在其任何点处都不可缩的空间，也存在在某一点处可缩、但并非在每一点处都可缩的空间（III, p. 321, 习题 1）。
\item
  设 E 是 n 维欧氏空间（或更一般地，R 上的拓扑向量空间），X 是 E 的一个子集。若对每个 \(y \in X\) 及每个 \(t \in I\)，点 \(tx + (1 - t)y\) 都属于 X，则称集合 X 在其一点 x 处是星形的。E 的凸子集（I, p. 122）在其每一点处都是星形的。
\end{enumerate}

在其一点 x 处星形的 E 的拓扑子空间 X 在该点处可缩。事实上，映射 \(\sigma\colon X \times I \to X\)（由 \(\sigma(y,t) = tx + (1 - t)y\) 定义）是在 x 处带基点的同伦，它把 \(Id_{X}\) 连接到以 \(\{x\}\) 为像的常值映射。

特别地，R 的每个区间、欧氏空间的每个凸子集，或更一般地，R 上拓扑向量空间的每个凸子集，都在其每一点处可缩。

\begin{enumerate}
\def\labelenumi{\arabic{enumi})}
\setcounter{enumi}{4}
\tightlist
\item
  我们将在后面（TA, V）证明，欧氏球面 \(S_n\)（\(n\geqslant1\)）不与一点同伦等价。可分无穷维希尔伯特空间的单位球面在其每一点处都可缩（EVT, V, p. 71, 习题 13）。
\end{enumerate}

\subsection{相对同伦}

设 X、Y 是拓扑空间，A 是 X 的一个子空间，B 是 Y 的一个子空间。

设 \(f:\mathrm{X}\to\mathrm{Y}\) 是满足 \(f(\mathrm{A})\subset\mathrm{B}\) 的连续映射。我们称 f 是偶 (X,A) 到偶 (Y,B) 上的同伦等价，

如果存在连续映射 \(g\colon Y\to X\)（满足 \(g(B)\subset A\)）、同伦 \(\sigma\colon X\times I\to X\)（在 A 上固定，把 \(\text{Id}_X\) 与 \(g\circ f\) 相联系）以及同伦 \(\tau\colon Y\times I\to Y\)（在 B 上固定，把 \(\text{Id}_Y\) 与 \(f\circ g\) 相联系）。

假设这些条件满足。那么：

\begin{enumerate}
\def\labelenumi{\alph{enumi})}
\item
  映射 f 是从 X 到 Y 的同伦等价，映射 g 是从 Y 到 X 的同伦等价，且这两个同伦等价在同伦意义下互为逆。
\item
  此时映射 g 是偶 (Y,B) 到偶 (X,A) 上的同伦等价，称为 f 在同伦意义下的一个逆。
\item
  映射 f 与 g 通过限制到子空间，诱导出从 A 到 B 上的互逆的同胚。
\end{enumerate}

若存在偶 (X,A) 到偶 (Y,B) 上的同伦等价，则称偶 (X,A) 与 (Y,B) 同伦等价。关系“X、Y 是拓扑空间，A 是 X 的子空间，B 是 Y 的子空间，且偶 (X,A) 与 (Y,B) 同伦等价”是一个等价关系。

\subsection{收缩映射与收缩}

定义 4 —— 设 X 是拓扑空间，A 是 X 的一个子集。

我们把 X 到 A 上的收缩映射定义为从 X 到 A 的一个连续映射，它是 A 到 X 中的典范单射的收缩映射。若存在 X 到 A 上的收缩映射，则称 X 可收缩到 A 上，或称 A 是 X 的一个收缩核。

设 X 是拓扑空间，A 是 X 的一个子空间。

设存在 X 到 A 上的一个收缩映射 \(r\)。子空间 A 等同于点 \(x \in \mathrm{X}\)（满足 \(x = r(x)\)）所成的集合。若 X 是分离空间，则 A 在 X 中是闭的。

根据下面的引理，子空间 A 是 X 的收缩核的充要条件是：从 A 到拓扑空间 Y 的每个连续映射都可扩张为从 X 到 Y 的连续映射。

引理 1 —— 设 X 与 Y 是拓扑空间，\(f\colon X\to Y\) 是连续映射。下列条件等价：

\begin{enumerate}
\def\labelenumi{(\roman{enumi})}
\item
  对每个拓扑空间 Z 及每个连续映射 \(g: X \to Z\)，存在连续映射 \(g': Y \to Z\)，使得 \(g = g' \circ f\)；
\item
  映射 f 是单射且有一个连续的收缩映射；
\item
  映射 f 定义一个从 X 到其像 \(f(\mathrm{X})\) 上的同胚，且该像是 Y 的一个收缩核。
\end{enumerate}

设 \(f\) 是单射，\(r: \mathrm{Y} \to \mathrm{X}\) 是映射 \(f\) 的一个连续收缩映射。对每个连续映射 \(g: \mathrm{X} \to \mathrm{Z}\)，映射 \(g' = g \circ r: \mathrm{Y} \to \mathrm{Z}\) 满足 \(g' \circ f = g \circ r \circ f = g\)。这就证明了 (ii)⇒(i)。

设条件 (i) 满足。令 \(g \colon X \to X\) 为恒同映射。由假设，存在连续映射 \(g' \colon Y \to X\)，使得 \(g' \circ f = g\)。由此推出映射 \(f\) 是单射，且 \(g'\) 是 \(f\) 的一个连续收缩映射。

性质 (ii) 与 (iii) 的等价性是显然的。

定义 5 —— 设 X 是拓扑空间，A 是 X 的一个子集。我们把 X 到 A 上的收缩定义为在 A 上固定的一个同伦 \(\sigma\colon X\times I\to X\)，其起点是 X 的恒同映射，其终点是 X 到 A 上的一个收缩映射。若存在 X 到 A 上的收缩，则称 X 可收缩到 A 上，或称 A 是 X 的一个收缩。

注 —— 换言之，X 到 A 上的收缩是具有下列性质的同伦 \(\sigma\colon X \times I \to X\)：

\begin{enumerate}
\def\labelenumi{(\roman{enumi})}
\item
  \(\sigma(x,0)=x\)（对所有 \(x\in X\)）；
\item
  \(\sigma(x,1)\in\mathrm{A}\)（对所有 \(x\in\mathrm{X}\)）；
\item
  \(\sigma(x,t)=x\)（对所有 \(x\in\mathbf{A}\) 及每个 \(t\in\mathbf{I}\)）。
\end{enumerate}

存在 X 到 A 上的收缩的充要条件是：A 到 X 中的典范单射是偶 (A,A) 到偶 (X,A) 上的同伦等价。

设 \(a\) 是 X 的一点。X 到子空间 \(\{a\}\) 上的收缩恰好就是在 \(a\) 处带基点、把映射 \(\mathrm{Id}_{\mathrm{X}}\) 连接到以 \(\{a\}\) 为像的常值映射的同伦。因此，X 可收缩到 \(\{a\}\) 上当且仅当 X 在 \(a\) 处可缩（III, p. 234, 例 3）。

设 X 是拓扑空间，A 是 X 的一个子集。设 \(\sigma\) 是 X 到 A 上的一个收缩。由 \(r(x) = \sigma(x,1)\) 定义的从 X 到 A 的映射 r 是 X 到 A 上的收缩映射，而 \(\sigma\) 是连接 \(\mathrm{Id}_{\mathrm{X}}\) 与 \(r\) 的同伦。于是，关系式 \(r \circ i = \mathrm{Id}_{\mathrm{A}}\) 与 \(i \circ r = r\) 表明映射 \(i\) 与 \(r\) 是同伦等价，且在同伦意义下互逆。

定义 6 —— 沿用前面的记号，称 \(\sigma\) 为强收缩，如果再有 \(r(\sigma(x,t))=r(x)\)（对所有 \(x\in X\) 及所有 \(t\in I\)）。

例 —— 设 X 是 \(B_{n}\) 中原点的余集。把 \(X \times I\) 映到 X 的、由 \((x, t) \mapsto ((1 - t) + t \frac{1}{\|x\|})x\) 给出的映射是 X 到 \(S_{n-1}\) 上的一个强收缩。与之相伴的 X 到 \(S_{n-1}\) 上的收缩映射是由 \(x \mapsto x / \|x\|\) 给出的映射。

引理 2 —— 设 X 是拓扑空间，U 是 X 的开子集，\(\sigma\) 是 X 到 \(X-U\) 上的一个强收缩。那么 \(\sigma(U\times I)\subset\overline{U}\)。特别地，\(\sigma(U\times\{1\})\) 包含在 U 的边界中。

设 \(x \in \mathrm{U}\)。实数 \(t \in \mathbf{I}\)（满足 \(\sigma(x, t) \in \mathrm{U}\)）所成的集合在 \(\mathbf{I}\) 中是开的且包含 0；设 \(s\) 是它的上确界。我们有 \(\sigma(x, s) \in \overline{\mathrm{U}}\)。若 \(s < 1\)，则 \(\sigma(x, s) \notin \mathrm{U}\)（由 \(s\) 的定义）；若 \(s = 1\)，由于 \(\sigma(x, 1) \in \mathrm{X} - \mathrm{U}\)，结论同样成立。于是 \(\sigma(x, s) \in \mathrm{Fr}(\mathrm{U})\)。由强收缩的定义，我们又有 \(\sigma(x, s) = \sigma(\sigma(x, s), 1) = \sigma(x, 1)\)；特别地，\(\sigma(x, 1) \in \overline{\mathrm{U}}\)。因此，\(\sigma(x, 1)\) 属于边界 \(\overline{\mathrm{U}} \cap (\mathrm{X} - \mathrm{U})\)（\(\mathrm{U}\) 的）。

设 \(t \in \mathbf{I}\)；若 \(\sigma(x,t) \notin \mathrm{U}\)，同样有 \(\sigma(x,t) = \sigma(\sigma(x,t),1) = \sigma(x,1)\)，因此 \(\sigma(x,t) \in \overline{\mathrm{U}}\)，引理得证。

\subsection{映射的柱}

设 X 与 Y 是拓扑空间，f 是从 X 到 Y 的连续映射。用 U 表示空间 \(U_{1} = X \times I\) 与空间 \(U_{2} = Y\) 的拓扑和空间，并经由典范单射把 \(U_{1}\) 与 \(U_{2}\) 等同于 U 的子空间。设 R 是 U 上最细的等价关系，使得 \((x,1)\)（\(U_{1}\) 的点）与 \(f(x)\)（\(U_{2}\) 的点）等价（对每个 \(x \in X\)）。

定义 7 —— 我们把商拓扑空间 U/R 称为映射 f 的柱，记作 \(\mathrm{Cyl}(f)\)。

用 \(\alpha_{f}\colon\mathrm{X}\times\mathbf{I}\to\mathrm{Cyl}(f)\) 与 \(\beta_{f}\colon\mathrm{Y}\to\mathrm{Cyl}(f)\) 表示从 U 到 U/R 上的典范满射在 \(\mathrm{U}_{1}\) 与 \(\mathrm{U}_{2}\) 上的限制。映射 \(\alpha_{f}\) 是一个同伦，其终点为 \(\beta_{f}\circ f\)。

命题 4（柱的泛性质）

设 Z 是拓扑空间，\(\beta: \mathrm{Y} \to \mathrm{Z}\) 是连续映射，\(\alpha: \mathrm{X} \times \mathbf{I} \to \mathrm{Z}\) 是终点为 \(\beta \circ f\) 的同伦。存在唯一的连续映射 \(\varphi\)（从 \(\operatorname{Cyl}(f)\) 到 Z），使得 \(\alpha = \varphi \circ \alpha_f\) 且 \(\beta = \varphi \circ \beta_f\)。

映射 \(\psi\)（从 U 到 Z，与 \(\alpha\) 在 \(\times\) 上、与 \(\beta\) 在 Y 上一致）是连续的。我们有 \(\alpha(x,1)=\beta(f(x))\)（对所有 \(x\in X\)），因此 \(\psi\) 过渡到商，给出连续映射 \(\varphi\)（从 \(\mathrm{Cyl}(f)\) 到 Z），使得 \(\alpha=\varphi\circ\alpha_{f}\) 且 \(\beta=\varphi\circ\beta_{f}\)。由于 \(\alpha_{f}\) 与 \(\beta_{f}\) 的像覆盖 \(\mathrm{Cyl}(f)\)，\(\varphi\) 是满足这些关系式的从 \(\mathrm{Cyl}(f)\) 到 Z 的唯一映射。

例 1 —— 设 X' 与 Y' 是拓扑空间，\(f': \mathrm{X}' \to \mathrm{Y}'\) 是连续映射。设 \(u: \mathrm{X} \to \mathrm{X}'\) 与 \(v: \mathrm{Y} \to \mathrm{Y}'\) 是满足 \(f' \circ u = v \circ f\) 的连续映射。存在唯一的连续映射 \(w: \operatorname{Cyl}(f) \to \operatorname{Cyl}(f')\)，使得 \(w \circ \alpha_f = \alpha_{f'} \circ (u \times \operatorname{Id}_\mathbf{I})\) 且 \(w \circ \beta_f = \beta_{f'} \circ v\)。

把命题 4 应用于 Z=Y 与 \(\beta = \operatorname{Id}_{Y}\)，存在唯一的连续映射 \(\gamma_{f}\colon\operatorname{Cyl}(f)\to\operatorname{Y}\)，使得有 \(\gamma_{f}(\alpha_{f}(x,s))=f(x)\) 与 \(\gamma_{f}(\beta_{f}(y))=y\)（对 \(x\in X\)、\(s\in I\) 及 \(y\in Y\)）。

命题 5 —— 映射 \(\alpha_{f}\) 诱导 X × {[}0,1{[} 到 \(\mathrm{Cyl}(f)\) 的一个开子集上的同胚。映射 \(\beta_{f}\) 定义 Y 到该开集的余集上的同胚。映射 \(\beta_{f} \circ \gamma_{f}\) 是 \(\mathrm{Cyl}(f)\) 到 \(\beta_{f}(\mathrm{Y})\) 上的连续收缩映射。

集合 X × {[}0,1{[} 是 U 的开子集，对关系 R 是饱和的，且 R 在 X × {[}0,1{[} 中诱导的等价关系是相等关系。由此即得第一个断言（TG, I, p. 23, 推论 10）。

\(\alpha_{f}(\mathrm{X}\times[0,1])\) 的余集是 \(\beta_{f}(\mathrm{Y})\)。由于 \(\gamma_{f}\circ\beta_{f}=\mathrm{Id}_{\mathrm{Y}}\)，\(\beta_{f}\) 定义 Y 到 \(\beta_{f}(\mathrm{Y})\) 上的同胚，且 \(\beta_{f}\circ\gamma_{f}\) 是 \(\beta_{f}(\mathrm{Y})\) 到 \(\mathrm{Cyl}(f)\) 中的典范单射的连续收缩映射。

\(\beta_{f}(\mathrm{Y})\)（\(\operatorname{Cyl}(f)\) 的闭子空间）称为 \(f\) 的柱的柱底。映射 \(\beta_{f} \circ \gamma_{f}\) 称为 \(\operatorname{Cyl}(f)\) 到其柱底上的典范收缩映射。

考虑由下式定义的映射 \(\sigma_{1}\colon\mathrm{U}_{1}\times\mathbf{I}\to\mathrm{Cyl}(f)\) 与 \(\sigma_{2}\colon\mathrm{U}_{2}\times\mathbf{I}\to\mathrm{Cyl}(f)\)：

\[
\begin{array}{c c} \sigma_ {1} ((x, s), t) = \alpha_ {f} (x, (1 - t) s + t) & \text {for} (x, s) \in \mathrm{X} \times \mathbf {I} \text {and} t \in \mathbf {I} \\ \sigma_ {2} (y, t) = \beta_ {f} (y) & \text {for} y \in \mathrm{Y} \text {and} t \in \mathbf {I}. \end{array}
\]

它们是连续的。对 \(x \in X\) 及 \(t \in I\)，我们有

\[
\sigma_ {1} ((x, 1), t) = \alpha_ {f} (x, 1) = \beta_ {f} (f (x)) = \sigma_ {2} (f (x), t).
\]

因此存在唯一的映射 \(\sigma_{f}\)（从 \(\mathrm{Cyl}(f)\times\mathbf{I}\) 到 \(\mathrm{Cyl}(f)\)），使得 \(\sigma_{f}\circ(\alpha_{f}\times\mathrm{Id}_{\mathbf{I}})=\sigma_{1}\) 且 \(\sigma_{f}\circ(\beta_{f}\times\mathrm{Id}_{\mathbf{I}})=\sigma_{2}\)。映射 \(\sigma_{f}\) 是连续的（I, p. 19, 命题 10）。

命题 6 —— 映射 \(\sigma_{f}\colon\mathrm{Cyl}(f)\times\mathbf{I}\to\mathrm{Cyl}(f)\) 是 \(\mathrm{Cyl}(f)\) 到其柱底上的一个强收缩。其终点是 \(\mathrm{Cyl}(f)\) 到其柱底上的典范收缩映射。

映射 \(\sigma_f\) 是一个同伦。关系式 \(\sigma_f(\alpha_f(x,s),0) = \alpha_f(x,s)\) 与 \(\sigma_f(\beta_f(y),0) = \beta_f(y)\) 表明 \(\sigma_f\) 的起点是 \(\mathrm{Cyl}(f)\) 的恒同映射。用 \(r_f\) 表示 \(\sigma_f\) 的终点。关系式

\[
\sigma_ {f} \left(\alpha_ {f} (x, s), 1\right) = \alpha_ {f} (x, 1) = \beta_ {f} (f (x)) = \left(\beta_ {f} \circ \gamma_ {f}\right) \left(\alpha_ {f} (x, s)\right)
\]

与 \(\sigma_{f}(\beta_{f}(y),1) = \beta_{f}(y) = (\beta_{f}\circ \gamma_{f})(\beta_{f}(y))\) 表明 \(r_f\) 是典范收缩映射 \(\beta_{f}\circ \gamma_{f}\)（\(\mathrm{Cyl}(f)\) 到其柱底上的）。

对 \((x,s)\in \mathbf{X}\times \mathbf{I}\) 及 \(t\in \mathbf{I}\)，我们有

\[
r _ {f} \left(\sigma_ {f} \left(\alpha_ {f} (x, s), t\right)\right) = r _ {f} \left(\alpha_ {f} (x, s (1 - t) + t)\right) = \beta_ {f} (f (x)) = r _ {f} \left(\alpha_ {f} (x, s)\right).
\]

对 \(y \in Y\) 及 \(t \in I\)，有 \(r_{f}(\sigma_{f}(\beta_{f}(y), t)) = r_{f}(\beta_{f}(y))\)。于是，\(\sigma_{f}\) 是 \(\operatorname{Cyl}(f)\) 到其柱底上的强收缩。

映射 \(\sigma_{f}\) 称为 \(f\) 的柱到其柱底上的典范收缩。

注 —— 1) 设 A 是 X × I 的子集，A₁ 是使 (x,1) ∈ A 的点 x ∈ X 所成的集合。那么我们有 \(\bar{\alpha}_{f}^{-1}(\alpha_{f}(A)) = A \cup f^{-1}(f(A_{1})) \times \{1\}\) 与 \(\beta_{f}^{-1}(\alpha_{f}(A)) = f(A_{1})\)。因此，若 f 是闭（相应地，开）的，则映射 \(\alpha_{f}\) 亦然。

\begin{enumerate}
\def\labelenumi{\arabic{enumi})}
\setcounter{enumi}{1}
\tightlist
\item
  设映射 \(f\) 是逆紧的。设 \(P\) 是 \(\operatorname{Cyl}(f)\) 的一点。若 \(P = \alpha_f(x,t)\)，其中 \(0 \leqslant t < 1\) 且 \(x \in X\)，则我们有 \(\overline{\alpha}_f^1(P) = (x,t)\)。在相反情形，存在 \(y \in Y\) 使得 \(P = \beta_f(y)\)
\end{enumerate}

且 \(\bar{\alpha}_{f}^{1}(P) = \bar{f}^{1}(y) \times \{1\}\)。这证明映射 \(\alpha_{f}\) 的纤维是拟紧的。由于 \(\alpha_{f}\) 是闭映射，故它是逆紧的（TG, I, p. 75, 定理 1）。

由命题 5 及 TG, I, p. 72, 命题 2，映射 \(\beta_f\) 本身也是逆紧的。因此，若 f 逆紧，则从 U 到 \(\mathrm{Cyl}(f)\) 上的典范满射是逆紧的，从而特别地是普遍严格的（I, p. 20, 推论）。

\begin{enumerate}
\def\labelenumi{\arabic{enumi})}
\setcounter{enumi}{2}
\tightlist
\item
  若空间 X 与 Y 是分离的，则 f 的柱也是分离的。事实上，设 z 与 \(z'\) 是 \(\mathrm{Cyl}(f)\) 的两个不同点，我们来证明它们有互不相交的邻域。我们区分三种情形。
\end{enumerate}

- 存在 \((x,t)\) 与 \((x',t') \in \mathrm{X} \times [0,1[\)，使得 \(z = \alpha_f(x,t)\) 且 \(z' = \alpha_f(x',t')\)。

在此情形，断言由如下事实得出：空间 \(X \times [0,1]\) 是分离的（TG, I, p. 54, 命题 7），且映射 \(\alpha_{f}\) 诱导该空间到 \(\operatorname{Cyl}(f)\) 的一个开子空间上的同胚。

- 存在 \((x,t)\in \mathrm{X}\times [0,1[\) 使得 \(z = \alpha_{f}(x,t)\)，且存在 \(y^{\prime}\in \mathrm{Y}\) 使得 \(z^{\prime} = \beta_{f}(y^{\prime})\)。

此时，\(\alpha_{f}(\mathrm{X}\times [0,\frac{t + 1}{2}])\) 与 \(\operatorname {Cyl}(f)\mathbf{-}\alpha_{f}(\mathrm{X}\times [0,\frac{t + 1}{2}])\) 是 \(z\) 与 \(z^{\prime}\) 在 \(\operatorname {Cyl}(f)\) 中的互不相交的开邻域。

- 存在 \(y\) 与 \(y'\in Y\)，使得 \(z=\beta_f(y)\) 且 \(z'=\beta_f(y')\)。

在此情形，\(y\neq y'\)；由于 Y 是分离的，存在 y 在 Y 中的开邻域 V 以及开邻域 \(V'\)（\(y'\) 在 Y 中的），使得 \(V\cap V'=\varnothing\)。于是，\((\beta_f\circ\gamma_f)^{-1}(V)\) 与 \((\beta_f\circ\gamma_f)^{-1}(V')\) 是 \(\operatorname{Cyl}(f)\) 的互不相交的开子集，分别包含 \(\beta_f(y)\) 与 \(\beta_f(y')\)。

\subsection{同伦扩张性质}

定义 8 —— 设 X 是拓扑空间，A 是 X 的一个子集。若对每个拓扑空间 Y、每个连续映射 \(f:\mathrm{X}\to\mathrm{Y}\) 以及每个同伦 \(\sigma:\mathrm{A}\times\mathbf{I}\to\mathrm{Y}\)（其终点为映射 \(f|_{\mathrm{A}}\)），都存在同伦 \(\tau:\mathrm{X}\times\mathbf{I}\to\mathrm{Y}\)（它扩张 \(\sigma\)，且终点为映射 f），则称偶 (X,A) 具有同伦扩张性质。

注 1 —— 设 X 是拓扑空间，A 是 X 的子集，且偶 (X,A) 具有同伦扩张性质。设 Y 是拓扑空间，\(f:\mathrm{X}\to\mathrm{Y}\) 是连续映射，\(\sigma:\mathrm{A}\times\mathrm{I}\to\mathrm{Y}\) 是起点为映射 \(f|_{\mathrm{A}}\) 的同伦。映射 \(\overline{\sigma}:\mathrm{A}\times\mathrm{I}\to\mathrm{Y}\)（由 \((a,t)\mapsto\sigma(a,1-t)\) 定义）是终点为 \(f|_{\mathrm{A}}\) 的同伦；设 \(\tau:\mathrm{X}\times\mathrm{I}\to\mathrm{Y}\) 是扩张 \(\overline{\sigma}\) 且终点为 f 的同伦。于是映射 \(\overline{\tau}:\mathrm{X}\times\mathrm{I}\to\mathrm{Y}\)（由 \((x,t)\mapsto\tau(x,1-t)\) 给出）是扩张 \(\sigma\) 且起点为映射 f 的同伦。

设 X 是拓扑空间，A 是 X 的子空间，\(i: A \to X\) 是典范单射。用 \(\alpha_{i}: A \times I \to \text{Cyl}(i)\) 与 \(\beta_{i}: X \to \text{Cyl}(i)\) 表示典范映射。设 \(j: \text{Cyl}(i) \to X \times I\) 是唯一的连续映射，使得有 \(j(\alpha_{i}(a, s)) = (i(a), s)\) 与 \(j(\beta_{i}(x)) = (x, 1)\)（对 \(a \in A\)、\(s \in I\) 及 \(x \in X\)）。它是单射；其像是子空间 \((\mathbf{A} \times \mathbf{I}) \cup (\mathbf{X} \times \{1\})\)（\(X \times I\) 的）。若 A 在 X 中闭，则映射 j 是闭的。它并不总是严格的（III, p. 325, 习题 17）。

命题 7 —— 沿用上述记号，下列断言等价：

\begin{enumerate}
\def\labelenumi{(\roman{enumi})}
\item
  偶 (X, A) 具有同伦扩张性质；
\item
  对每个拓扑空间 Y 及每个连续映射 \(g\colon \operatorname{Cyl}(i) \to \operatorname{Y}\)，存在连续映射 \(g'\colon X \times I \to Y\)，使得 \(g = g' \circ j\)；
\item
  单射 j 有一个连续的收缩映射；
\item
  映射 j 是严格的，且存在 X × I 到 j 的像上的一个收缩。
\end{enumerate}

设偶 (X,A) 具有同伦扩张性质。设 \(g:\operatorname{Cyl}(i)\to Y\) 是连续映射；令 \(\sigma=g\circ\alpha_i\)，\(f=g\circ\beta_i\)。映射 \(f:X\to Y\) 是连续的，而 \(\sigma:A\times I\to Y\) 是终点为映射 \(f|A\) 的同伦。因此存在同伦 \(g':X\times I\to Y\)（终点为 f，且扩张 \(\sigma\)）。我们有 \(g'\circ j(\alpha_i(a,s))=g'(a,s)=\sigma(a,s)=g(\alpha_i(a,s))\) 与 \(g'\circ j(\beta_i(x))=g'(x,1)=f(x)=g(\beta_i(x))\)（对 \(a\in A\)、\(s\in I\) 及 \(x\in X\)）。于是 \(g'\circ j=g\)，从而得到 (ii)。

反之，设断言 (ii) 成立，我们来证明偶 (X,A) 具有同伦扩张性质。设 Y 是拓扑空间，\(f:X\to Y\) 是连续映射，\(\sigma:A\times I\to Y\) 是终点等于 \(f|A\) 的同伦。存在唯一的连续映射 \(g:\mathrm{Cyl}(i)\to\mathrm{Y}\)，使得 \(g\circ\alpha_i=\sigma\) 且 \(g\circ\beta_i=f\)（III, p. 238, 命题 4）。于是，映射 \(g':\mathrm{X}\times\mathbf{I}\to\mathrm{Y}\) 若满足 \(g=g'\circ j\)，便是终点为 f 且扩张 \(\sigma\) 的同伦。

  断言 (ii) 与 (iii) 的等价性是 III, p. 235 引理 1 的特殊情形。

设单射 \(j\colon \mathrm{Cyl}(i)\to \mathrm{X}\times \mathbf{I}\) 有一个连续收缩映射 \(r\)。映射 \(j\) 定义从 \(\mathrm{Cyl}(i)\) 到子空间 \(\mathrm{T} = (\mathrm{A}\times \mathbf{I})\cup (\mathrm{X}\times \{1\})\)（\(\mathrm{X}\times \mathbf{I}\) 的）上的同胚（出处同上）。因此它是严格的。

令 \(\rho = \mathrm{pr}_1\circ j\circ r\)，\(\theta = \mathrm{pr}_2\circ j\circ r\)。若 \(x\in \mathrm{A}\) 或 \(s = 1\)，我们有 \(\rho (x,s) = x\) 与 \(\theta (x,s) = s\)。对 \((x,s)\in \mathrm{X}\times \mathbf{I}\) 及 \(t\in \mathbf{I}\)，令

\[
\sigma ((x, s), t) = (\rho (x, (1 - t) + s t), (1 - t) s + t \theta (x, s)).
\]

映射 \(\sigma\colon(\mathrm{X}\times\mathbf{I})\times\mathbf{I}\to\mathrm{X}\times\mathbf{I}\) 是连续的。对 \((x,s)\in\mathrm{X}\times\mathbf{I}\)，我们有

\[
\sigma ((x, s), 0) = (\rho (x, 1), s) = (x, s)
\]

以及

\[
\sigma ((x, s), 1) = (\rho (x, s), \theta (x, s)) = j \circ r (x, s);
\]

对 \(x\in \mathbf{X}\) 及 \(t\in \mathbf{I}\)，我们有

\[
\sigma ((x, 1), t) = (\rho (x, 1), (1 - t) + t \theta (x, 1)) = (x, 1),
\]

而对 \((x,s)\in\mathbf{A}\times\mathbf{I}\) 及 \(t\in\mathbf{I}\)，我们有

\[
\sigma ((x, s), t) = (x, (1 - t) s + t s) = (x, s).
\]

因此，\(\sigma\) 是 \(X\times I\) 到 T 上的一个收缩。这表明断言 (iii) 蕴涵断言 (iv)。

蕴涵关系 (iv)⇒(iii) 是显然的。

注 2 —— 设 X 是分离拓扑空间，A 是 X 的子集，且偶 (X,A) 具有同伦扩张性质。用 \(i\colon A\to X\) 表示典范单射。设 \(j\colon\operatorname{Cyl}(i)\to X\times\mathbf{I}\) 是典范单射，\(r\) 是映射 j 的一个连续收缩映射。子空间 \(j(\operatorname{Cyl}(i))\) 等于偶 \((x,t)\in X\times\mathbf{I}\)（满足 \(j(r(x,t))=(x,t)\)）所成的集合。由于空间 \(X\times\mathbf{I}\) 是分离的，子集 \(j(\operatorname{Cyl}(i))\) 在 \(X\times\mathbf{I}\) 中是闭的。于是集合 A——它等于 \(x\in X\)（使 \((x,0)\in j(\operatorname{Cyl}(i))\)）所成的集合——是 X 的闭子集。

引理 3 —— 设 X 与 Y 是拓扑空间，\(p \colon X \to Y\) 是逆紧且开的连续映射，\(f \colon X \to \mathbf{R}\) 是连续映射。映射 \(g \colon Y \to \overline{\mathbf{R}}\)（由 \(y \mapsto \sup_{x\in p^{-1}(y)} f(x)\) 给出）是连续的。

设 \(b \in \mathbf{Y}\)。由于 \(p\) 逆紧，其纤维 \(\overline{p}^1(b)\) 是拟紧空间；因此我们有 \(g(b) \in \mathbf{R} \cup \{-\infty\}\)。

设 \(m \in R\) 满足 \(g(b) < m\)。对所有 \(a \in \overline{p}^{1}(b)\)，我们有 \(f(a) \leqslant g(b) < m\)；设 \(V_{a}\) 是 a 在 X 中的邻域，使得 \(f(x) < m\)（对所有 \(x \in V_{a}\)）。诸集合 \(V_{a}\) 的并 V 是 \(\overline{p}^{1}(b)\) 在 X 中的邻域。由引理 5（I, p. 75），存在 b 在 Y 中的邻域 W，使得 \(\overline{p}^{1}(W) \subset V\)。对所有 \(y \in W\)，我们有 \(g(y) \leqslant m\)。这证明 g 在 b 处是上半连续的。

现在证明 \(g\) 在 \(b\) 处是下半连续的。我们可以假设 \(g(b)\in\mathbf{R}\)。设 \(m\in\mathbf{R}\) 满足 \(m<g(b)\)。设 \(a\in\overline{p}^1(b)\) 满足 \(m<f(b)\)；于是设 V 是 a 在 X 中的邻域，使得 \(f(x)>m\)（对所有 \(x\in V\)）。由此得出 \(g(y)>m\)（对所有 \(y\in p(V)\)）。由于 p 是开的，\(p(V)\) 是 b 的邻域，因此 g 是下半连续的。

引理至此证毕。

定理 1 —— 设 X 是拓扑空间，A 是 X 的闭子空间；用 i : A → X 表示典范单射。下列断言等价：

\begin{enumerate}
\def\labelenumi{(\roman{enumi})}
\item
  偶 (X, A) 具有同伦扩张性质；
\item
  存在连续映射 \(\varphi\colon\mathrm{X}\to\mathbf{I}\)，使得 \(\mathrm{A}=\overline{\varphi}^{1}(0)\)，并存在在 A 上固定、终点为 X 的恒同映射的同伦 \(\sigma\colon\mathrm{X}\times\mathbf{I}\to\mathrm{X}\)，使得有 \(\sigma(x,0)\in\mathrm{A}\)（对每个点 \(x\in\mathrm{X}\)，满足 \(\varphi(x)\neq1\)）；
\item
  存在连续映射 \(\varphi\colon\mathbf{X}\to\mathbf{R}_{+}\)，使得 \(\mathrm{A}=\overline{\varphi}^{1}(0)\)，并存在在 A 上固定的同伦 \(\sigma\colon\overline{\varphi}^{1}(\mathbf{I})\times\mathbf{I}\to\mathrm{X}\)，使得 \(\sigma(x,1)=x\) 且 \(\sigma(x,0)\in\mathrm{A}\)（对所有 \(x\in\overline{\varphi}^{1}(\mathbf{I})\)）；
\item
  存在连续映射 \(\varphi\colon X\to\mathbf{R}_{+}\)，使得 \(A\subset\overline{\varphi}^{1}(0)\)，并存在连续映射
\end{enumerate}

\[
\sigma \colon \{(x, t) \in \mathrm{X} \times \mathbf {I} \mid t + \varphi (x) \geqslant 1 \} \to \mathrm{X}
\]

使得有 \(\sigma(x,1) = x\)（对所有 \(x \in X\)）且 \(\sigma(x,1 - \varphi(x)) \in A\)（对所有 \(x \in X\)，满足 \(\varphi(x) \leqslant 1\)）。

我们用 i 表示 A 到 X 中的典范单射，用 \(\alpha_{i}\colon A\times I\to\operatorname{Cyl}(i)\) 与 \(\beta_{i}\colon X\to\operatorname{Cyl}(i)\) 表示典范映射，并用 j 表示从 \(\operatorname{Cyl}(i)\) 到 \(X\times I\) 的映射，它满足：\(j(\alpha_{i}(x,s))=(x,s)\)（对 \((x,s)\in A\times I\)），\(j(\beta_{i}(x))=(x,1)\)（对 \(x\in X\)）。

设偶 (X,A) 具有同伦扩张性质，并设 \(r:X\times I\to\operatorname{Cyl}(i)\) 是映射 j 的一个连续收缩映射（III, p. 241, 命题 7）。用 \(\sigma\) 表示连续映射 \(\operatorname{pr}_1\circ j\circ r\)（从 \(X\times I\) 到 X）。对每个 \((x,t)\in X\times I\)，我们有 \(|\operatorname{pr}_2(j(r(x,t)))-t|\leqslant1\)。由于 I 是紧的，第一投影 \(\operatorname{pr}_1:X\times I\to X\) 是逆紧映射；它也是开的。由引理 3，映射 \(\varphi:X\to I\)（由

\[
x \mapsto \sup _ {t \in \mathbf {I}} \left(\left| \operatorname{pr} _ {2} (j (r (x, t))) - t \right|\right)
\]

给出）是连续的。

对 \(x \in \mathrm{X}\)，我们有 \(\sigma(x,1) = x\)。设 \(x \in \mathrm{X}\) 满足 \(\sigma(x,0) \notin \mathrm{A}\)；则 \(\mathrm{pr}_2(j(r(x,0))) = 1\) 且 \(\varphi(x) \geqslant 1\)。对 \(x \in \mathrm{A}\) 及 \(t \in \mathbf{I}\)，我们有 \(j(r(x,t)) = (x,t)\)；因此 \(\sigma(x,t) = x\) 且 \(\varphi(x) = 0\)。反之，设 \(x \in \mathrm{X}\) 满足 \(\varphi(x) = 0\)。于是有 \(\mathrm{pr}_2(j(r(x,t)) \leqslant t\)（对所有 \(t \in \mathbf{I}\)）；若 \(t < 1\)，这蕴涵 \(j(r(x,t)) \in \mathrm{A} \times \mathbf{I}\)；由于 A 在 X 中闭，我们因此有 \(j(r(x,1)) \in \mathrm{A} \times \mathbf{I}\)，从而 \(x \in \mathrm{A}\)。

这表明 (i) 蕴涵 (ii)。

设 \(\varphi\) 与 \(\sigma\) 是满足断言 (ii) 中性质的映射。令 \(\varphi_1 = 2\varphi\)，并设 \(\sigma_1\) 是 \(\sigma\) 在 \(\overline{\varphi}_1^1 (\mathbf{I})\times \mathbf{I}\) 上的限制。我们有 \(\overline{\varphi}_1^1 (0) = \mathrm{A}\)；对 \(a\in \mathrm{A}\)，有 \(\sigma_1(a,t) = a\)（对所有 \(a\in \mathrm{A}\)）。设 \(x\in \mathrm{X}\) 满足 \(\varphi_{1}(x)\leqslant 1\)；我们有 \(\sigma_1(x,1) = x\)；此外，\(\varphi (x) = \varphi_{1}(x) / 2 < 1\)，因此 \(\sigma_{1}(x,0)\in \mathrm{A}\)。于是 (ii) 蕴涵 (iii)。

我们来证明 (iii) 蕴涵 (iv)。设 \(\varphi\colon X\to\mathbf{R}\) 与 \(\sigma\colon\overline{\varphi}^{1}(\mathbf{I})\times\mathbf{I}\to X\) 如命题所述；令 \(B=\overline{\varphi}^{1}(\mathbf{I})\)，\(C=\overline{\varphi}^{1}([1,+\infty[)\)。

设 \(u_1\) 是从 B × I 到 X 的映射：\(u_1(x,t) = \sigma(x,1 - (1 - t)/2\varphi(x))\)（若 \(t + 2\varphi(x) \geqslant 1\) 且 \(\varphi(x) > 0\)），否则 \(u_1(x,t) = \sigma(x,0)\)。由下面的引理 4，它是连续的。

设 \(u_2\) 是从 B × I 到 X 的映射：\(u_2(x,t)=\sigma(x,\sup(0,1-2(1-t)(1-\varphi(x))))\)。它是连续的。对 \(x\in\mathrm{B}\)（满足 \(\varphi(x)=1/2\)）及 \(t\in\mathbf{I}\)，我们有 \(u_1(x,t)=\sigma(x,t)=u_2(x,t)\)。用 \(u:\mathrm{B}\times\mathbf{I}\to\mathrm{X}\) 表示如下的映射：\(u(x,t)=u_1(x,t)\)（若 \(\varphi(x)\leqslant1/2\)），\(u(x,t)=u_2(x,t)\)（若 \(1/2<x\leqslant1\)）。它是连续的，因为它在 B×I 的闭子空间 \(\overline{\varphi}^{1}([0,1/2])\times\mathbf{I}\) 与 \(\overline{\varphi}^{1}([1/2,1])\times\mathbf{I}\) 上的限制都是连续的（TG, I, p. 19, 命题 4）。

对 \(x\in X\)（满足 \(\varphi(x)=1\)）及 \(t\in\mathbf{I}\)，我们有 \(u(x,t)=u_2(x,t)=\sigma(x,1)=x\)。因此存在唯一的映射 \(\tau:X\times I\to X\)，它与 u 在 \(B\times I\) 上一致，与映射 \(\mathrm{pr}_1\) 在 \(C\times I\) 上一致；它是连续的（出处同上）。

设 \(x \in X\)；可以验证 \(\tau(x,1) = 1\)。此外，若 \(x \in A\)，则 \(\varphi(x) = 0\)，因此 \(\tau(x,t) = u_{1}(x,t) = \sigma(x,0) = x\)。最后，若 \(2\varphi(x) \leqslant 1\)，则 \(\varphi(x) \leqslant 1/2\)，因此 \(\tau(x,1 - 2\varphi(x)) = \sigma(x,0) \in A\)。

这就证明了断言 (iv) 成立。

最后，设断言 (iv) 满足，我们来证明偶 (X, A) 具有同伦扩张性质。

用 \(C_{1}\)（相应地，\(C_{2}\)）表示偶 \((x,t)\in\mathbf{X}\times\mathbf{I}\)（满足 \(t+\varphi(x)\leqslant1\)（相应地，\(t+\varphi(x)\geqslant1\)））所成的集合。它们都是闭集。对 \((x,t)\in\mathrm{C}_{1}\)，我们有 \(\sigma(x,1-\varphi(x))\in\mathrm{A}\)；设 \(\rho_{1}\colon\mathrm{C}_{1}\to\mathrm{Cyl}(i)\) 是由 \((x,t)\mapsto\alpha_{i}(\sigma(x,1-\varphi(x)),t+\varphi(x))\) 给出的映射；它是连续的。再设 \(\rho_{2}\colon\mathrm{C}_{2}\to\mathrm{Cyl}(i)\) 是由 \((x,t)\mapsto\beta_{i}(\sigma(x,t))\) 给出的连续映射。对 \((x,t)\in\mathrm{C}_{1}\cap\mathrm{C}_{2}\)，我们有 \(t+\varphi(x)=1\)，因此

\[
\rho_ {1} (x, t) = \alpha_ {i} (\sigma (x, 1 - \varphi (x)), 1) = \beta_ {i} (\sigma (x, 1 - \varphi (x)) = \rho_ {2} (x, t).
\]

因此存在唯一的映射 \(\rho\colon X\times\mathbf{I}\to\mathrm{Cyl}(i)\)，它与 \(\rho_{1}\) 在 \(C_{1}\) 上一致，与 \(\rho_{2}\) 在 \(C_{2}\) 上一致；它是连续的（TG, I, p. 19, 命题 4）。

对 \(x \in A\) 及 \(t \in I\)，我们有 \(\varphi(x) = 0\)，因此 \(t + \varphi(x) \leqslant 1\)。

\[
\rho (j (\alpha_ {i} (x, t))) = \rho (x, t) = \alpha_ {i} (\sigma (x, 1), t) = \alpha_ {i} (x, t).
\]

对 \(x\in \mathrm{X}\)，我们还有

\[
\rho (j (\beta_ {i} (x))) = \rho (x, 1) = \beta_ {i} (\sigma (x, 1)) = \beta_ {i} (x).
\]

由于映射 \(\alpha_{i}\) 与 \(\beta_{i}\) 的像覆盖 \(\mathrm{Cyl}(i)\)，由此得出映射 \(\rho\) 是映射 j 的连续收缩映射。于是，偶 (X, A) 具有同伦扩张性质（III, p. 241, 命题 7），定理证毕。

引理 4 —— 设 X 与 Y 是拓扑空间，\(\varphi\colon\mathrm{X}\to\mathbf{R}_{+}\) 是连续映射，并令 \(\mathrm{A}=\overline{\varphi}^{1}(0)\)。设 \(\sigma\colon\mathrm{X}\times\mathbf{I}\to\mathrm{Y}\) 是在 A 上固定的同伦。映射 \(\sigma'\colon\mathrm{X}\times\mathbf{I}\to\mathrm{Y}\) 把 \((x,s)\) 映到 \(\sigma(x,s/\varphi(x))\)（若 \(s<\varphi(x)\)）且映到 \(\sigma(x,1)\)（若 \(s\geqslant\varphi(x)\)），它是连续的。

映射 \(\sigma'\) 在闭子集 \(\overline{\varphi}^{1}([1, +\infty[)\times\mathbf{I}\) 的每一点处都连续；因此只需证明它在 \(\overline{\varphi}^{1}(\mathbf{I})\times\mathbf{I}\) 上的限制是连续的。于是我们可以假设 \(\varphi(\mathrm{X})\subset\mathbf{I}\)。设 C 与 C' 分别是由偶 \((x,s)\)（满足 \(s\leqslant\varphi(x)\) 与 \(s\geqslant\varphi(x)\)）组成的 X × I 的子空间。它们是闭的且覆盖 X × I。映射 \(\sigma'\) 在 C' 上连续；我们来证明它在 C 上连续。

设 \(\alpha\colon\mathbf{X}\times\mathbf{I}\to\mathbf{X}\times\mathbf{I}\) 是由 \(\alpha(x,t)=(x,t\varphi(x))\) 给出的连续映射。它的像等于 C，且 \(\sigma'\circ\alpha\) 是连续映射 \(\sigma\)。我们来证明 \(\alpha\) 是逆紧映射。事实上，考虑超滤子 \(\mathfrak{U}\)（在 \(\mathbf{X}\times\mathbf{I}\) 上）以及点 \((x,t)\in\mathbf{X}\times\mathbf{I}\)（附着于超滤子基 \(\alpha(\mathfrak{U})\)）。由于 \(\mathrm{pr}_{1}\circ\alpha=\mathrm{pr}_{1}\)，X 上的超滤子基 \(\mathrm{pr}_{1}(\mathfrak{U})\) 收敛于 \(x\)。由于 I 是紧的，存在点 \(s\in\mathbf{I}\) 使超滤子基 \(\mathrm{pr}_{2}(\mathfrak{U})\) 收敛于 \(s\)。于是 \(\mathfrak{U}\) 收敛于 \((x,s)\)。由于 \(\alpha\) 连续，超滤子基 \(\alpha(\mathfrak{U})\) 收敛于 \((x,s\varphi(x))\)。由于 I 是分离的，我们有 \(s\varphi(x)=t\)，于是 \(\alpha(x,s)=(x,t)\)，从而 \(\alpha\) 是逆紧的（TG, I, p. 75, 定理 1）。再由 I, p. 18, 例 2 及命题 9 推出 \(\sigma'\mid\mathrm{C}\) 是连续的。

推论 1 —— 设 X 是正规拓扑空间，A 是 X 的闭子空间。假设存在 A 在 X 中的邻域 V 及 V 到 A 上的一个收缩，以及连续映射 \(f: \mathrm{X} \to \mathbf{R}\)（满足 \(f^{-1}(0) = \mathrm{A}\)）。那么偶 (X,A) 具有同伦扩张性质。

设 \(\rho\colon V \times I \to V\) 是 \(V\) 到 \(A\) 上的一个收缩。由正规空间的定义（TG, IX, p. 41, 定义 1），存在连续映射 \(g\colon X \to I\)，它在 \(A\) 上取值 0，在 \(X - V\) 的每点处取值 1。设 \(\varphi\colon X \to R\) 是由 \(\varphi(x) = |f(x)| + g(x)\)（\(x \in X\)）给出的映射；它是连续的。我们有 \(\overline{\varphi}^{1}(0) = A\) 与 \(\overline{\varphi}^{1}(I) \subset V\)。设 \(\sigma\) 是从 \(\overline{\varphi}^{1}(I) \times I\) 到 \(X\) 的映射，由 \(\sigma(x,t) = \rho(x,1-t)\) 给出（对 \(x \in \overline{\varphi}^{1}(I)\) 及 \(t \in I\)）。对 \(x \in \overline{\varphi}^{1}(I)\)，我们有 \(\sigma(x,1) = \rho(x,0) = x\) 与 \(\sigma(x,0) = \rho(x,1) \in A\)。

映射 \(\varphi\) 与 \(\sigma\) 满足 III, p. 243 定理 1 断言 (iii) 的条件；因此，偶 (X,A) 具有同伦扩张性质。

例 1 —— 取球 \(\mathbf{B}_{n}\) 作为空间 X，球面 \(\mathbf{S}_{n-1}\) 作为子空间 A。若 \(V = X - \{0\}\)，则存在 V 到 \(\mathbf{S}_{n-1}\) 上的一个强收缩（III, p. 237, 例）。因此，偶 \((\mathbf{B}_{n}, \mathbf{S}_{n-1})\) 具有同伦扩张性质。

推论 2 —— 设 X 与 Y 是拓扑空间，A 是 X 的闭子空间，B 是 Y 的闭子空间。若偶 (X,A) 与 (Y,B) 都具有同伦扩张性质，则偶 (X × Y, (X × B) ∪ (A × Y)) 也具有同伦扩张性质。

设 \(\varphi\colon X\to\mathbf{R}_{+}\) 与 \(\sigma\colon\overline{\varphi}^{1}(\mathbf{I})\times\mathbf{I}\to X\)，以及 \(\varphi^{\prime}\colon Y\to\mathbf{R}_{+}\) 与 \(\sigma^{\prime}\colon\overline{\psi}^{1}(\mathbf{I})\times\mathbf{I}\to Y\)，分别是关于偶 (X,A) 与偶 (Y,B) 满足定理 1 断言 (iv) 的条件的映射。设 \(\psi\colon X\times Y\to\mathbf{R}_{+}\) 是由 \((x,y)\mapsto\inf(\varphi(x),\varphi^{\prime}(x))\) 给出的映射；它是连续的；我们还有 \(\psi(x,y)=0\)（对所有 \((x,y)\in X\times Y\)，满足 \(x\in A\) 或 \(y\in B\)）。对 \((x,y,t)\in X\times Y\times I\)（满足 \(t+\psi(x)\geqslant1\)），我们有 \(t+\varphi(x)\geqslant1\) 与 \(t+\varphi^{\prime}(x)\geqslant1\)。于是可以定义连续映射

\[
\tau \colon \{(x, y, t) \in \mathrm{X} \times \mathrm{Y} \times \mathbf {I} \mid t + \psi (x) \geqslant 1 \} \rightarrow \mathrm{X} \times \mathrm{Y}
\]

如下：令 \(\tau(x,y,t) = (\sigma(x,t),\sigma'(y,t))\)。对所有 \((x,y)\in X\times Y\)，我们有 \(\tau(x,y,1) = (\sigma(x,1),\sigma'(y,1)) = (x,y)\)。此外，若 \(\psi(x,y)\leqslant 1\)，则 \(\varphi(x)\leqslant 1\) 或 \(\varphi'(y)\leqslant 1\)，从而 \(\sigma(x,1)\in A\) 或 \(\sigma'(y,1)\in B\)；这蕴涵 \(\tau(x,y,1 - \psi(x))\) 属于 \((A\times Y)\cup (X\times B)\)。这就验证了定理 1 的断言 (iv)，推论得证。

例 2 —— *下面是偶 (X, A) 具有同伦扩张性质的其他情形。

\begin{enumerate}
\def\labelenumi{(\roman{enumi})}
\item
  空间 X 是 C \(^{1}\) 类仿紧微分流形，A 是 X 的闭子流形。鉴于推论 \(^{1}\)，这由如下事实得出：X 是完全正规的（IX, p. 103, 习题 11），且 A 在 X 中有一个管状邻域；
\item
  空间 X 是胞腔复形，A 是相对于 X 的一个给定胞腔分解的子复形。*
\end{enumerate}
